\chapter{Recommender Systems}

\section{Introduzione e Concetti Fondamentali}

\subsection{Inquadramento e Contesto}
Negli ultimi decenni, la diffusione pervasiva del Web e delle piattaforme digitali ha portato a una \textbf{crescita esponenziale della quantità di informazioni, prodotti e servizi accessibili agli utenti}. Tale fenomeno prende il nome di \textbf{information overload} (sovraccarico informativo). \textbf{In molti contesti decisionali}, la presenza di un numero elevatissimo di opzioni/alternative (dette comunemente \textit{item}) rende la scelta autonoma da parte dell'utente un compito estremamente complesso, lungo e cognitivamente oneroso.



\textbf{Quindi il testo ci sta dicendo che nel Web, ci sono davvero tante informazioni, e prendere decisioni con cosi tante informazioni, può essere un compito complesso e estramamente oneroso.}



In tale scenario si inserisce la branca dell'Intelligenza Artificiale nota come \textbf{AI-Supported Decision Making}, il cui obiettivo è sviluppare sistemi (\textbf{tipo i sistemi di raccomandazione}) capaci di assistere l'utente nell'adozione di decisioni ottimali in contesti ad alta complessità.

\begin{definition}{Recommender System}
Un \textbf{Sistema di Raccomandazione} (\textit{Recommender System}, RS) è \textbf{un'applicazione software appartenente all'area dell'AI-supported decision making} che affronta l'information overload (quindi delle tantissime informazioni presenti nel web) mediante tecniche di \textbf{information filtering}. Il sistema seleziona e propone all'utente un \textbf{sottoinsieme ridotto di opzioni (\textit{shortlist}) ritenute di suo potenziale interesse, apprendendone preventivamente le preferenze su un determinato dominio} (ad es. film, musica, prodotti e-commerce, articoli di notizie).
\end{definition}

Da un punto di vista concettuale, un sistema di raccomandazione si colloca all'intersezione tra tre macro-aree:
\begin{enumerate}
    \item \textbf{Decision Making}: la teoria per supportare l'utente nell'esplorazione dello spazio delle alternative e nella scelta finale, cioè praticamente quello che il sistema decide di far vedere all'utente dopo la fase di filtraggio, ad esempio i primi 5-10 video di YouTube con piu hype.
    \item \textbf{Learning (Machine Learning \& Statistica)}: l'insieme di algoritmi necessari per estrarre pattern, modellare il comportamento degli utenti e predirne i gusti futuri a partire da dati storici o interazioni.
    \item \textbf{Information Filtering}: l'insieme di tecniche per filtrare il flusso informativo e ridurre il catalogo di partenza a pochi elementi rilevanti.
    Da un catalogo gigantesco con milioni di oggetti (film, prodotti, articoli),il filtro rimuove tutto ciò che non è rilevante per te, riducendo l'enorme flusso informativo a un set ristretto di poche decine di opzioni.
\end{enumerate}

Immagina Youtube:
\begin{itemize}
    \item \textbf{Information Filtering}: Dagli oltre un miliardo di video presenti sul server, il filtro ne seleziona una centinaia pertinenti.   
    \item \textbf{Learning}: Gli algoritmi di ML stimano la probabilità che tu clicchi e guardi ciascuno di quei 100 video basandosi sulla tua cronologia.   
    \item \textbf{Decision Making}: La home page ti presenta i primi 5-10 video organizzati con miniature accattivanti e titoli chiari per renderti semplice la scelta di cosa guardare ad un primo colpo d'occhio.
\end{itemize}

\subsection{Contesti Applicativi e Benefici}
I sistemi di raccomandazione costituiscono oggi il nucleo operativo dei principali servizi aziendali e piattaforme web:
\begin{itemize}
    \item \textbf{E-commerce} (es. Amazon, Alibaba): suggerimento di prodotti correlati o acquisti frequenti associati.
    \item \textbf{Social Networks} (es. Facebook, Instagram, LinkedIn): raccomandazione di nuovi contatti, contenuti multimediali, gruppi o inserzioni pubblicitarie.
    \item \textbf{Entertainment \& Streaming} (es. Netflix, Spotify, YouTube): raccomandazione di film, brani musicali o video personalizzati.
\end{itemize}

L'adozione di un sistema di raccomandazione genera benefici tangibili sia per l'utente finale sia per il provider del servizio:

\paragraph{Benefici per l'Utente:}
\begin{itemize}
    \item \textbf{Restringimento dello spazio delle scelte}: riduzione drastica del tempo necessario per trovare un contenuto d'interesse.
    \item \textbf{Selezione degli item più rilevanti}: identificazione di prodotti ad alta utilità individuale.
    \item \textbf{Supporto alla conoscenza del dominio}: ausilio nella comprensione dell'offerta disponibile all'interno del catalogo.
    \item \textbf{Esplorazione di nuove alternative}: scoperta di prodotti o contenuti che l'utente non avrebbe cercato autonomamente.
\end{itemize}

\paragraph{Benefici per il Provider:}
\begin{itemize}
    \item \textbf{Aumento della Customer Satisfaction e Retention}: un'esperienza d'uso personalizzata incrementa la fidelizzazione dell'utente.
    \item \textbf{Promozione del catalogo (\textit{Item Promotion})}: possibilità di guidare l'attenzione degli utenti verso specifici prodotti (es. contenuti di catalogo secondario o a maggior margine economico).
    \item \textbf{Acquisizione di informazioni sull'utente}: raccolta continua di dati di interazione utile per profilazioni di mercato avanzate.
\end{itemize}


\section{Architettura e Tassonomia dei Dati}
\subsection{Flusso Input/Output di un Recommender System}

Dal punto di vista architetturale, un \textit{Recommender System} (sistema
di raccomandazione) può essere rappresentato come una componente che riceve
in ingresso diverse tipologie di informazioni, le elabora mediante uno o
più algoritmi e produce in uscita un insieme ordinato di elementi
(\textit{item}) ritenuti rilevanti per uno specifico utente.

L'obiettivo del sistema è individuare, all'interno di un catalogo, gli
elementi che presentano la maggiore probabilità di soddisfare gli
interessi, le preferenze o le esigenze dell'utente target. A tal fine,
il sistema può sfruttare sia informazioni relative agli elementi
disponibili e all'utente, sia dati derivanti dalle interazioni della
comunità e conoscenze strutturate sul dominio applicativo.

In generale, le principali categorie di input sono quattro:
\textit{Items Data}, \textit{User Data}, \textit{Community Data} e
\textit{Domain Model}.

\paragraph{Items Data.}
Gli \textit{Items Data} comprendono le informazioni che descrivono gli
elementi presenti nel catalogo del sistema di raccomandazione. Tali
informazioni possono includere caratteristiche descrittive, attributi,
categorie di appartenenza e altre proprietà utili a identificare e
confrontare gli elementi.

Ad esempio, in un sistema che raccomanda film, gli \textit{Items Data}
possono comprendere il titolo, il genere, la durata, l'anno di
produzione, il regista e gli attori di ciascun film. In un sistema
di raccomandazione di prodotti, possono invece includere il prezzo,
la marca, la categoria merceologica e le caratteristiche tecniche
dei prodotti disponibili.

Queste informazioni permettono al sistema di rappresentare gli
elementi del catalogo e di individuare quelli che presentano
caratteristiche compatibili con le preferenze dell'utente.

\paragraph{User Data.}
Gli \textit{User Data} rappresentano le informazioni relative
all'utente target, ovvero la persona alla quale il sistema deve
proporre le raccomandazioni. Possono comprendere preferenze esplicite,
interessi dichiarati, esigenze specifiche e, quando disponibili,
informazioni ricavate dal comportamento individuale.

Ad esempio, in un sistema di raccomandazione di film, un utente
potrebbe manifestare una preferenza per i generi fantascientifico
e thriller, oppure indicare di preferire film di durata inferiore
a due ore. Il sistema può utilizzare tali informazioni per
selezionare gli elementi maggiormente coerenti con il profilo
dell'utente.

Le informazioni individuali non devono necessariamente essere
fornite direttamente dall'utente: alcune preferenze possono essere
inferite dalle sue interazioni precedenti, come visualizzazioni,
acquisti o valutazioni. In tal caso, il sistema utilizza le evidenze
disponibili per costruire una rappresentazione degli interessi
individuali.

\paragraph{Community Data.}
I \textit{Community Data} comprendono le informazioni generate
dall'insieme degli utenti attraverso le loro interazioni con gli
elementi del catalogo. Possono includere valutazioni espresse,
recensioni, acquisti, visualizzazioni, clic e altre forme di
feedback esplicito o implicito.

Ad esempio, in un sistema di raccomandazione di film, questi dati
possono indicare quali film sono stati apprezzati dai diversi utenti
e quali titoli sono stati maggiormente visualizzati. Il sistema può
utilizzare tali informazioni per individuare elementi popolari
oppure per confrontare i comportamenti di utenti con preferenze
simili.

Un'applicazione particolarmente importante dei \textit{Community Data}
è il \textit{Collaborative Filtering} (filtraggio collaborativo),
una famiglia di tecniche che sfrutta le interazioni della comunità
per produrre raccomandazioni. Ad esempio, se due utenti hanno
espresso valutazioni simili su diversi film, il sistema può
utilizzare tale somiglianza per suggerire a uno dei due film
apprezzati dall'altro.

È importante distinguere questi dati dagli \textit{User Data}:
mentre questi ultimi riguardano specificamente l'utente target,
i \textit{Community Data} descrivono le interazioni della comunità
nel suo complesso e possono essere utilizzati per ricavare
informazioni utili alla personalizzazione.

\paragraph{Domain Model.}
Il \textit{Domain Model} è una rappresentazione strutturata della
conoscenza relativa al dominio applicativo in cui opera il sistema
di raccomandazione. Può comprendere concetti, relazioni semantiche,
regole e vincoli che descrivono il significato degli elementi e
le modalità con cui essi sono collegati tra loro.

A differenza degli \textit{Items Data}, che descrivono le proprietà
dei singoli elementi del catalogo, il \textit{Domain Model} permette
di rappresentare la struttura concettuale del dominio e le relazioni
che intercorrono tra i suoi elementi. Non si limita quindi a
memorizzare informazioni sui singoli oggetti, ma organizza la
conoscenza necessaria per interpretarli nel loro contesto.

Ad esempio, nel dominio cinematografico, gli \textit{Items Data}
possono indicare che un determinato film appartiene al genere
fantascientifico ed è diretto da un certo regista. Il
\textit{Domain Model} può invece rappresentare relazioni generali
quali l'appartenenza di un film a uno o più generi, la relazione
tra un film e il suo regista oppure i collegamenti semantici tra
generi e temi narrativi.

Tale conoscenza può essere sfruttata per individuare relazioni
significative tra elementi diversi. Se un utente apprezza un film
di fantascienza, ad esempio, il sistema può prendere in
considerazione altri film dello stesso genere o caratterizzati
da temi simili. La condivisione di una caratteristica non implica
necessariamente che il nuovo elemento sarà apprezzato dall'utente,
ma costituisce un'informazione che può contribuire alla
valutazione della sua rilevanza.

Un ulteriore esempio riguarda la raccomandazione di materiale
didattico: un modello del dominio può rappresentare le relazioni
tra argomenti, corsi ed esami, nonché eventuali conoscenze
preliminari richieste. In questo modo, il sistema può proporre
risorse coerenti non soltanto con gli interessi dello studente,
ma anche con il percorso di apprendimento da seguire.

Il \textit{Domain Model} non coincide necessariamente con un
algoritmo di raccomandazione: costituisce piuttosto una fonte
di conoscenza strutturata che il sistema può utilizzare insieme
alle altre informazioni disponibili.

\paragraph{Dipendenza dal dominio applicativo.}
Le quattro categorie di input costituiscono una classificazione
generale, applicabile a diversi contesti. Il loro contenuto
concreto dipende tuttavia dal dominio applicativo considerato.
Ad esempio, in un sistema cinematografico gli elementi del
catalogo sono film, mentre in un sistema di e-commerce sono
prodotti e in un sistema musicale sono canzoni o album.

Analogamente, il \textit{Domain Model} deve rappresentare i
concetti e le relazioni specifici del contesto considerato.
La stessa architettura generale può quindi essere riutilizzata
in domini differenti, adattando i dati, le conoscenze e le
tecniche di raccomandazione alle caratteristiche del problema.

\paragraph{Output del sistema.}
A partire dalle informazioni ricevute, il \textit{Recommender
System} elabora le informazioni disponibili e produce in uscita
una \textit{shortlist} di elementi rilevanti per l'utente target.
La shortlist è generalmente un sottoinsieme del catalogo,
selezionato in base alla rilevanza stimata dei singoli elementi
e, quando previsto, ordinato dal più al meno pertinente.

Ad esempio, un catalogo cinematografico può contenere migliaia
di film, ma il sistema ne seleziona soltanto alcuni da proporre
a un determinato utente. La selezione può dipendere dalle
caratteristiche dei film, dalle preferenze individuali, dalle
interazioni della comunità e dalle relazioni rappresentate nel
modello del dominio.

Il flusso generale può essere sintetizzato come segue:

\begin{equation}
\begin{aligned}
\text{Input} = \{&
\text{Items Data},\,
\text{User Data},\\
&\text{Community Data},\,
\text{Domain Model}
\}\\[4pt]
&\downarrow\\[-2pt]
&\boxed{\text{Recommender System}}\\[-2pt]
&\downarrow\\[2pt]
&\text{Output: Shortlist di Item Rilevanti}
\end{aligned}
\end{equation}

Di seguito viene illustrata la rappresentazione dettagliata delle quattro tipologie di dati d'ingresso per un Recommender System operante in un dominio cinematografico.

\paragraph{Esempio Integrato di Input JSON.}
Considerando una piattaforma di streaming video, l'input complessivo fornito al motore di raccomandazione per elaborare un suggerimento personalizzato per un utente target (\texttt{user\_42}) è strutturato come segue:

\begin{lstlisting}[language=json]
{
  "items_data": [
    {
      "item_id": "movie_101",
      "titolo": "Inception",
      "anno": 2010,
      "durata_minuti": 148,
      "regista": "Christopher Nolan",
      "generi": ["Fantascienza", "Thriller"],
      "attori_principali": ["Leonardo DiCaprio", "Joseph Gordon-Levitt"]
    },
    {
      "item_id": "movie_102",
      "titolo": "Interstellar",
      "anno": 2014,
      "durata_minuti": 169,
      "regista": "Christopher Nolan",
      "generi": ["Fantascienza", "Drammatico"],
      "attori_principali": ["Matthew McConaughey", "Anne Hathaway"]
    }
  ],
  "user_data": {
    "user_id": "user_42",
    "preferenze_esplicite": {
      "generi_preferiti": ["Fantascienza", "Thriller"],
      "durata_massima_minuti": 180
    },
    "storico_individuale_inferito": {
      "film_visti_recentemente": ["movie_101"],
      "tempo_completamento_medio": "95%"
    }
  },
  "community_data": [
    {
      "user_id": "user_07",
      "item_id": "movie_101",
      "rating": 5,
      "interazione": "like"
    },
    {
      "user_id": "user_07",
      "item_id": "movie_102",
      "rating": 5,
      "interazione": "like"
    }
  ],
  "domain_model": {
    "relazioni_generi": {
      "Fantascienza": {
        "generi_correlati": ["Cyberpunk", "Space Opera"],
        "temi_chiave": ["Viaggi nel tempo", "Spazio"]
      }
    },
    "regole_semantiche": [
      {
        "regola": "Affinita_Regista_Genere",
        "descrizione": "I film del medesimo regista appartenenti allo stesso genere presentano un'elevata rilevanza incrociata."
      }
    ]
  }
}
\end{lstlisting}

\paragraph{Elaborazione dell'Input.}
Il sistema elabora questi dati in ingresso combinando i vincoli dell'utente (\emph{User Data}) con la descrizione dei film (\emph{Items Data}). Attraverso i \emph{Community Data}, riconosce la somiglianza tra \texttt{user\_42} e \texttt{user\_07} (Collaborative Filtering), mentre mediante il \emph{Domain Model} interpreta la correlazione tematica e di regia tra i titoli, generando infine la raccomandazione del film \emph{Interstellar}.

\subsection{Modellazione degli Item e Domain Model}

\subsubsection{Items Data}
La modalità con cui gli item vengono rappresentati all'interno del sistema dipende dal dominio e dalla tecnica algoritmica adottata:
\begin{itemize}
    \item \textbf{Dati Non Strutturati}: l'item è identificato unicamente da un'etichetta alfanumerica o ID privo di attributi descrittivi espliciti. Questo approccio è tipico del \textit{Collaborative Filtering}, che opera esclusivamente sulle matrici di interazione senza necessitare del contenuto degli item.
    \item \textbf{Dati Strutturati}: l'item è descritto da un identificativo e da un insieme strutturato di attributi (o \textit{feature}) che ne specificano le caratteristiche fisiche, tecniche o concettuali (ad esempio, per uno smartphone: RAM, processore, megapixel della fotocamera, prezzo).
    \item \textbf{Descrizioni Testuali}: testo in linguaggio naturale che descrive l'item (es. sinossi di un libro o trama di un film).
\end{itemize}

I sistemi \textit{Content-Based} e i sistemi basati su \textit{Bayesian Preference Elicitation} richiedono obbligatoriamente dati strutturati per effettuare il confronto tra le preferenze e gli attributi.

\subsubsection{Domain Model}
Il \textit{Domain Model} ha lo scopo di modellare la conoscenza del dominio per arricchire la descrizione degli item e definire relazioni esplicite tra di essi. Include:
\begin{itemize}
    \item \textbf{Schemi di attributi}: definizioni formali dei domini e dei tipi di dato degli attributi.
    \item \textbf{Tassonomie}: organizzazioni gerarchiche di categorie e sottocategorie (es. \textit{Film} $\rightarrow$ \textit{Documentario} $\rightarrow$ \textit{Documentario Storico}).
    \item \textbf{Ontologie}: grafi di conoscenza ed esplicitazioni formali di concetti e relazioni (Linked Data, DBpedia).
    \item \textbf{Metriche di similitudine}: funzioni matematiche in grado di quantificare quanto due item siano vicini nello spazio delle caratteristiche.
\end{itemize}

\subsection{Modellazione dell'Utente e della Comunità}

\subsubsection{User Data}
Rappresenta il profilo del singolo utente:
\begin{itemize}
    \item \textbf{Informazioni personali/demografiche}: età, sesso, titolo di studio, professione, localizzazione geografica.
    \item \textbf{Informazioni di utilizzo}: frequenza d'accesso, grafo sociale, follower/following.
    \item \textbf{Modello delle preferenze}: rappresentazione esplicita degli interessi, pesi attribuiti ai vari attributi o vera e propria \textit{funzione di utilità}.
\end{itemize}

\subsubsection{Community Data}
Registra i comportamenti collettivi di tutti gli utenti registrati:
\begin{itemize}
    \item \textbf{Interazioni User-Item}: le valutazioni storiche (es. punteggi, stelle, "like"), i brani ascoltati, gli articoli acquistati.
    \item \textbf{Statistiche sugli item}: indicatori di popolarità globale (es. prodotti più venduti, film più visti).
    \item \textbf{Statistiche sugli utenti}: identificazione di utenti centrali, esperti o \textit{influencer} della community.
\end{itemize}



\subsection{Tecniche di Acquisizione dei Dati: Esplicita vs Implicita}
Le informazioni sulle preferenze degli utenti possono essere raccolte mediante due modalità principali:

\subsubsection{Acquisizione Esplicita}
L'utente fornisce consapevolmente un giudizio o un'informazione al sistema.
\begin{itemize}
    \item \textbf{Esempi}: inserimento dati anagrafici, rating numerici (1--5 stelle), recensioni testuali, inserimento di "like/dislike", compilazione di questionari sulle preferenze.
    \item \textbf{Vantaggi (+)}: dati facili da interpretare semanticamente; elevata consapevolezza dell'utente sulle informazioni rilasciate.
    \item \textbf{Svantaggi (-)}: richiede sforzo attivo ed energia da parte dell'utente (alto \textit{carico cognitivo}); le attività di compilazione possono risultare noiose; i dati ricevuti possono contenere rumore o valutazioni incostanti.
\end{itemize}

\subsubsection{Acquisizione Implicita}
Il sistema osserva passivamente il comportamento dell'utente durante la normale fruizione del servizio.
\begin{itemize}
    \item \textbf{Esempi}: storico degli acquisti effettuati, tracciamento dei click, tempo di permanenza (\textit{dwell time}) su una pagina o su un video, cronologia delle ricerche o dei film visti.
    \item \textbf{Vantaggi (+)}: modalità totalmente non invasiva (non richiede alcuna azione aggiuntiva rispetto alla navigazione standard); garantisce una grande mole di dati continui.
    \item \textbf{Svantaggi (-)}: dati intrinsecamente ambigui (ad es. l'aver guardato un film per intero implica necessariamente che sia piaciuto? O l'utente potrebbe essersi addormentato?); mancanza di feedback negativi espliciti; potenziale percezione di intrusione nella privacy da parte dell'utente ("sensazione di essere spiati").
\end{itemize}


\section{Tipologia dei Sistemi di Raccomandazione}

I sistemi di raccomandazione si suddividono in due macro-categorie principali: \textbf{Non Personalizzati} e \textbf{Personalizzati}.

\subsection{Sistemi Non Personalizzati vs Personalizzati}

\paragraph{Sistemi Non Personalizzati (\textit{Non-Personalized}):}
Forniscono la medesima lista di raccomandazioni a tutti gli utenti (o a interi sottoinsiemi non profilati), basandosi esclusivamente su indicatori statistici aggregati del comportamento collettivo o su proprietà intrinseche degli item.
\begin{itemize}
    \item \textbf{Vantaggi}: applicabili immediatamente ai nuovi utenti privi di storico; costi computazionali contenuti grazie alla generatività omogenea degli output.
    \item \textbf{Svantaggi}: raccomandazioni poco accurate per le specifiche preferenze del singolo individuo.
\end{itemize}

\paragraph{Sistemi Personalizzati (\textit{Personalized}):}
Analizzano le interazioni pregresse del singolo utente per inferirne il profilo gusti e generare suggerimenti ad-hoc.
\begin{itemize}
    \item \textbf{Vantaggi}: raccomandazioni ad alta accuratezza e personalizzazione.
    \item \textbf{Svantaggi}: elevato carico computazionale (ogni utente richiede un'analisi dedicata); incapacità di operare quando un utente è privo di interazioni passate (\textbf{New User Cold-Start Problem}).
\end{itemize}

\subsection{Tassonomia dei Metodi Non Personalizzati}
I sistemi non personalizzati si strutturano principalmente in:
\begin{enumerate}
    \item \textbf{Popularity-based}: si basano sul calcolo di metriche aggregate sulla comunità.
    \begin{itemize}
        \item \textit{Top-ranked}: item con i punteggi medi più alti.
        \item \textit{Most viewed / Most clicked}: item con la massima frequenza di fruizione.
    \end{itemize}
    \item \textbf{Feature-based}: si basano su criteri descrittivi o promozionali gestiti dal provider.
    \begin{itemize}
        \item Item aggiunti di recente o in uscita imminente dal catalogo.
        \item Prodotti in sconto o promozionali.
        \item Categorie specifiche (es. "Serie e film storici").
        \item Item con riconoscimenti di critica.
    \end{itemize}
\end{enumerate}

\subsection{Tassonomia dei Metodi Personalizzati}

I sistemi di raccomandazione personalizzati si dividono in due grandi famiglie metodologiche in base alla natura dei dati utilizzati e al meccanismo di generazione dei suggerimenti:

\begin{enumerate}
    \item \textbf{Content-based / Knowledge-based:} raccomandano item che presentano caratteristiche e attributi simili a quelli espressamente richiesti dall'utente o da lui apprezzati in passato.
    \begin{itemize}
        \item \textbf{History-based:} analizzano lo storico delle interazioni passate dell'utente per modellarne il profilo. Se un utente fruisce frequentemente di contenuti con determinati attributi (es. film appartenenti al genere \textit{thriller}), il sistema proporrà altri item condividenti i medesimi metadati.
        \begin{itemize}
            \item \textit{Sfida principale:} definire ed elaborare adeguate metriche di similarità matematica (es. Cosine Similarity, TF-IDF) tra item rappresentati tramite vettori di attributi.
        \end{itemize}
        
        \item \textbf{Query-based:} sfruttano un linguaggio esplicito di interrogazione mediante il quale l'utente specifica requisiti rigidi (\textit{hard constraints}, es. prezzo massimo) o flessibili (\textit{soft constraints}, es. vicinanza al centro).
        \begin{itemize}
            \item \textit{Sfida principale:} progettare un \textit{domain model} sufficientemente ricco dal punto di vista semantico e un'interfaccia utente espressiva ma di semplice utilizzo.
        \end{itemize}
    \end{itemize}

    \item \textbf{Il Collaborative Filtering} (CF) è una tecnica di raccomandazione che si basa sull'idea che utenti con gusti simili tendono ad apprezzare gli stessi item, oppure che item apprezzati dagli stessi utenti possono essere correlati.

    La caratteristica principale è che il sistema sfrutta le interazioni degli utenti con gli item, come voti, valutazioni (rating) e acquisti, senza aver necessariamente bisogno di conoscere le caratteristiche degli item stessi.

    \paragraph{Esempio} :  Consideriamo un recommender system di film:

    - L'utente A ha apprezzato Interstellar e Inception.
    - L'utente B ha apprezzato gli stessi film, ma anche The Prestige.
    - Il sistema potrebbe raccomandare The Prestige all'utente A, sfruttando la somiglianza tra le interazioni dei due utenti.

    Il sistema non deve necessariamente sapere che questi film sono di fantascienza o che sono diretti da Christopher Nolan: può formulare la raccomandazione basandosi sulle valutazioni degli utenti.

    Il Collaborative Filtering si divide principalmente in due categorie:
    \begin{itemize}
        \item Neighborhood-based (User-based / Item-based): individua somiglianze tra utenti oppure tra item, utilizzandole direttamente per formulare le raccomandazioni.
        \item Model-based: apprende un modello a partire dalle interazioni disponibili, per individuare schemi nei dati e generare raccomandazioni.
    \end{itemize}

    In sintesi, la distinzione fondamentale è che il Collaborative Filtering sfrutta il comportamento collettivo degli utenti, anziché basarsi necessariamente sulle caratteristiche intrinseche degli item, come genere, descrizione o categoria.
\end{enumerate}


\section{Obiettivi di Valutazione di un Recommender System}

La progettazione e la scelta di un algoritmo di raccomandazione devono bilanciare molteplici obiettivi, sia di natura \textbf{operativa} (orientati all'esperienza utente) che di natura \textbf{strutturale} (orientati all'architettura software).

\subsection{Obiettivi Operativi}
\begin{itemize}
    \item \textbf{Relevance (Rilevanza)}: misura il grado di gradimento generale degli item raccomandati. Si valuta tracciando il tasso di selezione, i rating rilasciati o il tempo medio di utilizzo continuo del servizio.
    \item \textbf{Novelty (Novità)}: misura la capacità del sistema di suggerire item pertinenti che l'utente non conosceva o non avrebbe scoperto autonomamente.
    \item \textbf{Serendipity (Serendipità)}: misura la capacità di raccomandare item inaspettati e sorprendenti, estranei alle abitudini consolidate dell'utente ma ugualmente apprezzati (ad es., suggerire un ristorante afgano a un utente che ha sempre frequentato trattorie locali). A differenza della \textit{novelty}, la \textit{serendipity} richiede una forte componente di sorpresa positiva.
    \item \textbf{Diversity (Diversità)}: valuta la varietà degli item presenti nella shortlist proposta, evitando di raccomandare opzioni quasi identiche (es. evitare di proporre tutti i capitoli della stessa saga cinematografica in una singola lista).
\end{itemize}

\subsection{Obiettivi Strutturali e Sfide Architetturali}
\begin{itemize}
    \item \textbf{Scalability (Scalabilità)}: capacità dell'algoritmo di operare in modo efficiente su grandi volumi di dati ($10^7$--$10^8$ utenti e $10^4$--$10^6$ item).
    \begin{itemize}
        \item I metodi \textit{Non-Personalized} garantiscono la massima scalabilità.
        \item I metodi \textit{Neighborhood-based} soffrono di elevata complessità al crescere di utenti e item.
        \item I metodi basati su \textit{Preference Elicitation} richiedono elaborazioni statistiche in tempo reale che possono limitarne la scalabilità su larghissima scala.
    \end{itemize}

    \item \textbf{Explainability (Spiegabilità)}: la capacità del sistema di fornire una motivazione comprensibile all'utente riguardo alla raccomandazione effettuata.
    \begin{itemize}
        \item Metodi complessi basati su Machine Learning o vettori latenti agiscono spesso come "scatole nere" (\textit{black-box}) difficili da spiegare.
        \item I metodi \textit{Content-based} guidati da ontologie o attributi espliciti offrono le migliori basi per generare spiegazioni intuitive.
    \end{itemize}

    \item \textbf{Cold-Start Problem (Problema dell'Avvio a Freddo)}: si verifica quando un nuovo utente (\textit{New User Cold-Start}) o un nuovo item (\textit{New Item Cold-Start}) fa il suo ingresso nel sistema privo di uno storico di interazioni.
    \begin{itemize}
        \item \textit{Mitigazione per gli utenti}: ricorso iniziale a raccomandazioni non personalizzate o somministrazione di un questionario d'ingresso (\textit{onboarding}).
        \item \textit{Mitigazione per gli item}: impiego di approcci content-based che analizzano le proprietà dell'item indipendentemente dai rating ricevuti.
    \end{itemize}
\end{itemize}


\newpage


\section{User Rating Matrix (URM) e Tipologie di Rating}

\subsection{La Matrice User-Rating (URM)}
La grande maggioranza delle tecniche di raccomandazione (specialmente nel \textit{Collaborative Filtering}) formalizza l'insieme dei dati storici della community attraverso una struttura matriciale incompleta nota come \textbf{User Rating Matrix} (URM), indicata con $R$.

\begin{definition}{User Rating Matrix}
Sia $U = \{u_1, u_2, \dots, u_N\}$ l'insieme degli $N$ utenti e sia $I = \{i_1, i_2, \dots, i_M\}$ l'insieme degli $M$ item. La \textbf{User Rating Matrix} $R \in \mathbb{R}^{N \times M}$ è una matrice le cui righe corrispondono agli utenti e le cui colonne corrispondono agli item. L'elemento generico $r_{ui}$ rappresenta la valutazione espressa dall'utente $u$ nei confronti dell'item $i$:
\[
R = \begin{pmatrix}
r_{11} & r_{12} & \dots & r_{1M} \\
r_{21} & r_{22} & \dots & r_{2M} \\
\vdots & \vdots & \ddots & \vdots \\
r_{N1} & r_{N2} & \dots & r_{NM}
\end{pmatrix}
\]
Questa matrice contiene tutti gli item relativi ad un determinato dominio/catalogo.
\end{definition}


Poiché ogni singolo utente valuta tipicamente solo una frazione trascurabile dell'intero catalogo, la matrice $R$ risulta \textbf{fortemente sparsa} (\textit{sparse matrix}), ovvero la stragrande maggioranza delle sue celle è priva di valore (valore non osservato).

\subsection{Tipologie di Rating}
La tipologia di rating determina quali valori possono sostituire i numeri reali presenti nella matrice.
\begin{enumerate}
    \item \textbf{Rating Continui : } I valori possono essere qualsiasi \textbf{numero reale} all'interno di un intervallo.
    \begin{equation}
        r_{ui} \in [0,10]
    \end{equation}
    Esempio di matrice:
    \[
        R = \begin{pmatrix}
        8.7 & 6.25  \\
        9.1 & 4.83
        \end{pmatrix}
    \]
    \item \textbf{Rating Discreti : } I valori appartengono a un \textbf{INSIEME finito di valori ammessi}.
    \begin{equation}
        r_{ui} \in \{1,2,3,4,5\}
    \end{equation}
    Esempio di matrice:
    \[
        R = \begin{pmatrix}
        5 & 3  \\
        4 & 1
        \end{pmatrix}
    \]
    l'utente non può assegnare 3.7 stelle se la scala ammette solo questi cinque valori.
    \item \textbf{Rating Binari : } Sono ammessi soltanto due valori, che rappresentano due stati contrapposti.
    \begin{equation}
        r_{ui} \in \{0,1\}
    \end{equation}
    Esempio di matrice:
    \[
        R = \begin{pmatrix}
        1 & 0  \\
        1 & 1
        \end{pmatrix}
    \]
    Per esempio, 1 significa «Mi piace» e 0 significa «Non mi piace»
    \item \textbf{Rating Unari} : Si registra soltanto il feedback positivo.
    Esempio :  $r_{ui} = 1$ indica che l'utente ha cliccato su un prodotto.
    \[
        R = \begin{pmatrix}
        1 & 0 & 0  \\
        1 & 1 & 0
        \end{pmatrix}
    \]
    In questo esempio, 1 indica \textbf{un'interazione osservata}, mentre 0 significa che non è stata osservata alcuna interazione.
    \textbf{MA Lo zero non significa necessariamente che l'utente non gradisca l'item, semplicemente potrebbe non averlo mai visto}
    Contiene sono un informazione, mentre i rating binari, ne contengono due di informazioni, cioè ad esempio se all'utente è piaciuto o meno, mentre nei rating unari, conserviamo
    una singola informazione
\end{enumerate}

Sussiste un trade-off architetturale intrinseco tra \textbf{carico cognitivo dell'utente} e \textbf{potere discriminativo/risoluzione della scala}:
\begin{center}
\small
\textbf{Minimo Carico Cognitivo} (Unario $\rightarrow$ Binario $\rightarrow$ Discreto $\rightarrow$ Continuo) \textbf{Massimo Potere Discriminativo}
\end{center}

\subsection{Frequenza dei Rating e la Proprietà ``Long Tail''}

Analizzando la distribuzione empirica delle valutazioni all'interno di una User Rating Matrix (URM), si osserva frequentemente che gli item non ricevono tutti lo stesso numero di rating. In particolare, pochi item sono estremamente popolari, mentre un numero elevato di item riceve pochissime valutazioni. Questo fenomeno è noto come \textit{Long Tail Property} (proprietà della coda lunga).

\begin{definition}{Proprietà Long Tail}
La proprietà \textit{Long Tail} descrive una distribuzione fortemente sbilanciata della popolarità degli item: un numero esiguo di item raccoglie una grande quantità di interazioni, mentre una moltitudine di item riceve poche interazioni individuali.
\end{definition}

\paragraph{Short Head e Long Tail}

La distribuzione può essere suddivisa concettualmente in due parti:

\begin{itemize}
\item \textbf{Short Head (testa della distribuzione):} comprende i pochi item molto popolari, che raccolgono la maggior parte dei rating complessivi della comunità.
\item \textbf{Long Tail (coda della distribuzione):} comprende i numerosi item di nicchia, ciascuno caratterizzato da un numero ridotto di rating.
\end{itemize}

\paragraph{Esempio mediante una User Rating Matrix}

Consideriamo una matrice di valutazione $R$ con $N=5$ utenti e $M=6$ film. Supponiamo, per semplicità, di utilizzare una rappresentazione binaria delle interazioni, nella quale $1$ indica un'interazione osservata e $-$ indica l'assenza di una valutazione osservata.

$$
R =
\begin{array}{c|cccccc}
 & F_1 & F_2 & F_3 & F_4 & F_5 & F_6 \\
\hline
U_1 & 1 & 1 & 1 & - & - & - \\
U_2 & 1 & 1 & - & - & - & - \\
U_3 & 1 & 1 & - & - & - & - \\
U_4 & 1 & - & - & 1 & - & - \\
U_5 & 1 & - & - & - & 1 & -
\end{array}
$$

Per individuare la popolarità di ciascun item, è necessario contare le interazioni osservate nella corrispondente colonna. Definiamo quindi la popolarità dell'item $i$ come:

$$
p_i = \sum_{u=1}^{N}
\mathbf{1}\bigl[r_{ui}\text{ è osservato}\bigr]
$$

dove $\mathbf{1}[\cdot]$ è la funzione indicatrice, che assume valore $1$ se la condizione è verificata e $0$ altrimenti.

Applicando questa definizione alla matrice precedente, otteniamo:

$$
\begin{array}{c|cccccc}
\text{Item} & F_1 & F_2 & F_3 & F_4 & F_5 & F_6 \\
\hline
\text{Numero di rating} & 5 & 3 & 1 & 1 & 1 & 0
\end{array}
$$

Si osserva che $F_1$ ha ricevuto cinque interazioni, mentre $F_2$ ne ha ricevute tre. Questi due film costituiscono, nell'esempio, la \textit{Short Head}. Al contrario, $F_3$, $F_4$ e $F_5$ hanno ricevuto una sola interazione ciascuno e rappresentano item di nicchia riconducibili alla \textit{Long Tail}. Infine, $F_6$ non presenta interazioni osservate.

L'esempio è puramente illustrativo: una matrice così piccola consente di comprendere il fenomeno, ma non è sufficiente per dimostrare empiricamente la presenza di una distribuzione a legge di potenza (\textit{power law}).


\paragraph{Implicazioni per i Sistemi di Raccomandazione}

La proprietà Long Tail comporta alcune conseguenze importanti:

\begin{enumerate}
\item \textbf{Difficoltà di stima:} gli item con pochi rating forniscono informazioni limitate sulle preferenze degli utenti. Di conseguenza, le stime delle loro valutazioni o della loro rilevanza possono essere meno affidabili. Ad esempio, una valutazione media di $5$ stelle ottenuta da un solo utente è generalmente meno informativa di una media di $4.2$ stelle calcolata su migliaia di valutazioni.

\item \textbf{Popularity Bias:} molti algoritmi tendono a raccomandare gli item già popolari, anche perché dispongono di un maggior numero di interazioni osservate per apprenderne le caratteristiche e le preferenze associate. Si può così instaurare un circolo vizioso: gli item popolari vengono raccomandati più spesso, ricevono nuove interazioni e rafforzano ulteriormente la propria popolarità, mentre gli item di nicchia rischiano di essere trascurati.

\item \textbf{Impatto economico:} nei sistemi di e-commerce e streaming, raccomandare efficacemente gli item della Long Tail può favorire la scoperta di prodotti o contenuti meno conosciuti, valorizzare un catalogo più ampio e soddisfare interessi specifici degli utenti.

\end{enumerate}

\begin{warning}{Distinzione importante}
La popolarità di un item non coincide necessariamente con la sua qualità o con la valutazione media ricevuta. Un item molto popolare ha accumulato molte interazioni, ma ciò non implica che sia apprezzato da tutti gli utenti. Analogamente, un item con poche interazioni non è necessariamente di bassa qualità: potrebbe semplicemente essere poco conosciuto o rivolgersi a un pubblico di nicchia.
\end{warning}

In sintesi, la proprietà Long Tail riguarda la distribuzione del numero di interazioni tra gli item, non i singoli valori della matrice $R$. Per studiarla, si calcola la popolarità di ciascun item contando le valutazioni osservate nella relativa colonna e si analizza successivamente la distribuzione dei conteggi ottenuti.


\paragraph{Implicazioni per i Sistemi di Raccomandazione:}
\begin{enumerate}
    \item \textbf{Difficoltà di stima}: la scarsità numerica di rating sugli item nella coda rende ardua l'estrazione di stime statisticamente attendibili per essi.
    \item \textbf{Popularity Bias}: molti algoritmi di raccomandazione tendono a prediligere e suggerire sistematicamente i pochi item popolari situati nella testa, ignorando la coda.
    \item \textbf{Impatto Economico}: gli item popolari sono spesso quelli su cui i provider hanno margini di profitto inferiori (o costi di licenza più alti). La raccomandazione efficace degli item di nicchia collocati nella \textit{Long Tail} rappresenta una delle principali leve di valore economico per le piattaforme e-commerce e streaming.
\end{enumerate}


\section{Algoritmi Non Personalizzati e Normalizzazione}

Analizziamo ora in dettaglio il funzionamento matematico delle principali tecniche non personalizzate basate sulla URM e le metodologie per la rimozione dei bias sistematici nelle valutazioni.

\subsection{Approccio Most Popular}
L'approccio \textbf{Most Popular} rappresenta una delle forme più semplici ed essenziali di raccomandazione non personalizzata. Il sistema calcola il numero totale di interazioni o valutazioni ricevute da ciascun item e restituisce in output gli item con il conteggio più elevato.

\begin{definition}{Punteggio Most Popular}
Sia $S_j = \{u \in U \mid r_{uj} \text{ è osservato}\}$ l'insieme degli utenti che hanno valutato l'item $j$. Il punteggio associato all'item $j$ è dato dal modulo di tale insieme:
\[
\text{Score}_{MP}(j) = |S_j| = N_j
\]
\textbf{Fissato un item, si scorrono le colonne, cioè gli utenti e si conservano gli utenti che hanno ratato l'item}
\end{definition}

\paragraph{Fondamento psicologico e limiti:}
L'approccio si basa sull'\textit{istinto gregario} (gli utenti tendono a conformarsi alle scelte collettive della massa). Tuttavia, quando applicato su una URM, questo metodo presenta un importante limite strutturale: \textbf{considera unicamente il numero di giudizi ricevuti, ignorando del tutto il valore numerico del rating espressi}. Un item con 100 voti negativi (es. rating $1/5$) risulterà posizionato più in alto rispetto a un item con 90 voti eccellenti (es. rating $5/5$).

\subsection{Approccio Top Ranked}
Per superare i limiti del conteggio puro, l'approccio \textbf{Top Ranked} ordina gli item in base alla media aritmetica dei rating osservati nella URM.

\begin{definition}{Punteggio Top Ranked}
Sia $N_j = |S_j|$ il numero di valutazioni osservate per l'item $j$. La stima del punteggio medio per l'item $j$ è definita come:
\[
\hat{r}_j = \frac{1}{N_j} \sum_{u \in S_j} r_{uj}
\]
Cioè praticamente, quello che si fa è scorrere sugli utenti che hanno assegnato una valutazione all'item, e per ognuno prendo la valutazione dell'item (che so che esisterà perchè la prendo dall'insieme $S_j$), si sommano le valutazioni e poi si dividono per il numero di valutazioni totali che ha avuto quell'item $j$
\end{definition}

\begin{example}{Fallimento della Media Semplice}
Si consideri una scala di rating discreta $1$--$5$. Siano dati due item $i$ e $j$ presenti nella URM:
\begin{itemize}
    \item L'item $i$ ha ricevuto $5$ valutazioni: $\{4, 3, 4, 4, 5\}$.
    \item L'item $j$ ha ricevuto $1$ sola valutazione: $\{5\}$.
\end{itemize}

Calcolando la media aritmetica semplice per entrambi gli item si ottiene:
\[
\hat{r}_i = \frac{4 + 3 + 4 + 4 + 5}{5} = \frac{20}{5} = 4.0
\]
\[
\hat{r}_j = \frac{5}{1} = 5.0
\]

In base all'approccio Top Ranked puro, l'item $j$ viene classificato al di sopra dell'item $i$ ($\hat{r}_j > \hat{r}_i$), poiché presenta una media pari a $5.0$. Tuttavia, la stima su $j$ si basa su un singolo voto, risultando totalmente priva di significatività e stabilità statistica. In contesti reali (es. TripAdvisor), questo fenomeno può favorire comportamenti collusivi o recensioni artificiali.
\end{example}

\subsection{Regolarizzazione mediante Shrink Term (Damped Mean)}
Per mitigare la distorsione causata da campioni di ampiezza ridotta, si introduce un fattore di regolarizzazione ad-hoc a denominatore, detto \textbf{Shrink Term} (o termine di smorzamento), indicato con $C > 0$.

\begin{definition}{Media Regolarizzata con Shrink Term}
La media regolarizzata di un item $h$ è definita come:
\[
\hat{r}_h = \frac{\sum_{k \in S_h} r_{kh}}{N_h + C}
\]
dove $N_h$ è il numero di valutazioni dell'item $h$ e $C$ è una costante \textbf{positiva} di regolarizzazione.
\end{definition}

\paragraph{Analisi dell'effetto dello Shrink Term:}
Lo shrink term agisce penalizzando severamente gli item che possiedono un numero ridotto di valutazioni mentre l'impatto di $C$ diventa trascurabile per gli item molto popolari.

Questo perchè matematicamente, quando $N_j$ è piccolo, aggiungendo una costante positiva al denominatore, si penalizza tutto il risultato, mentre se $N_j$ è già grande, aggiungere una costante $C$ cambia poco

\begin{example}{Applicazione dello Shrink Term}
Riconsideriamo l'esempio precedente applicando uno *shrink term* pari a $C = 1$:
\begin{itemize}
    \item Per l'item $i$ ($N_i = 5$, somma voti $= 20$):
    \[
    \hat{r}_i = \frac{4 + 3 + 4 + 4 + 5}{5 + 1} = \frac{20}{6} \approx 3.33
    \]
    \item Per l'item $j$ ($N_j = 1$, somma voti $= 5$):
    \[
    \hat{r}_j = \frac{5}{1 + 1} = \frac{5}{2} = 2.50
    \]
\end{itemize}
A seguito della regolarizzazione, l'item $i$ ottiene un punteggio superiore rispetto a $j$ ($\hat{r}_i = 3.33 > \hat{r}_j = 2.50$). Lo *shrink term* ha penalizzato correttamente l'item $j$ riducendone la stima a causa della scarsa numerosità campionaria.
\end{example}

\subsection{Limiti dei Metodi Most Popular e Top Ranked}
Sebbene rappresentino un baseline semplice e veloce, i metodi non personalizzati presentano rilevanti limitazioni operazionali:
\begin{enumerate}
    \item \textbf{Penalizzazione di Novelty e Serendipity}: gli item popolari rimangono bloccati nelle prime posizioni, impedendo agli utenti di scoprire contenuti di nicchia.
    \item \textbf{Feedback Loop Positivo (\textit{Rich-get-richer effect})}: il sistema consiglia continuamente prodotti già popolari, incrementandone ulteriormente le visualizzazioni e le valutazioni a scapito del resto del catalogo.
    \item \textbf{Concentrazione del catalogo}: insoddisfazione per i fornitori/produttori che vedono azzerata la visibilità dei prodotti di nicchia.
\end{enumerate}

\subsection{Normalizzazione della URM: Rimozione dei Bias Globali}
I rating espressi all'interno di una URM dipendono non solo dalla reale qualità degli item, ma anche da \textbf{bias sistematici} legati alle diverse abitudini degli utenti (utenti severi vs. utenti generosi) e \textbf{alla popolarità generale del dominio}.

Per uniformare i dati ed eliminare tali distorsioni, si applica una procedura di \textbf{normalizzazione della URM in 6 passaggi sequenziali}:

\subsubsection{Procedura Algoritmica di Normalizzazione}

\paragraph{Step 1: Calcolo della Media Globale ($\hat{g}$)}
Si calcola il valore medio di tutti i rating osservati nella matrice $R$:
\[
\hat{g} = \frac{\sum_{(i,j) \in S} r_{ij}}{|S|}
\]
dove $S = \{(i,j) \mid r_{ij} \text{ è osservato}\}$ e $|S|$ rappresenta il numero totale di celle non vuote nella URM. La media globale $\hat{g}$ rappresenta il livello basale delle valutazioni dell'intero sistema.

\paragraph{Step 2: Sottrazione della Media Globale}
Si calcola la matrice dei rating centrati rispetto alla media globale:
\[
r'_{ij} = r_{ij} - \hat{g} \quad \forall (i,j) \in S
\]

Se questa differenza è maggiore di 0, allora quel rating viene \textbf{considerato sopra la media globale}

Questo passaggio rimuove la tendenza generale dell'intero sistema (ad es. in molti contesti gli utenti lasciano una recensione prevalentemente quando un item è piaciuto, portando la media globale $\hat{g}$ verso valori medio-alti).

\textbf{Si lavora su una nuova matrice $R'$, dove nella posizione $(i,j)$ della matrice originale, se la posizione era vuota, rimane vuota anche in questa}

\paragraph{Step 3: Calcolo dell'Item Bias ($b_j$)}
L'\textbf{Item Bias} $b_j$ è una media. Ci dice \textbf{in media} quanto il rating dell'utente $i$ sull'item $j$ si discosta dalla baseline globale (media globale).Applicando un fattore di regolarizzazione $C$, si definisce:
\[
b_j = \frac{\sum_{i: (i,j) \in S} r'_{ij}}{N_j + C}
\]
dove $N_j$ è il numero totale di valutazioni ricevute dall'item $j$.

\textbf{Operativamente} quello che si fa è :
\begin{equation}
    b_j = r_{1j}-\hat{g} + r_{2j}-\hat{g} + \dots + r_{Nj} - \hat{g}
\end{equation}
Quindi $b_j$ misura in media, quanto $r_{ij}$ si discosti da $\hat{g}$

\begin{itemize}
    \item se $b_j$ è positivo, significa che la maggior parte delle volte $r_{ij}$ > $\hat{g}$, questo significa che le recensioni degli utenti su quel prodotto $j$, erano sopra la baseline globale.
    Quindi il prodotto è valutato positivamente.
\end{itemize}



\paragraph{Step 4: Sottrazione dell'Item Bias}
Quindi se $b_j$ = 0.4, questo significa che le recensioni degli utenti, si scostano dalla baseline globale in media di 0.4, cioè le $r_{ij}$ fissato $i$, si discostano dalla baseline globale $\hat{g}$ di 0.4
\[
r''_{ij} = r'_{ij} - b_j = r_{ij} - \hat{g} - b_j \quad \forall (i,j) \in S
\]
Questo significa che faccio:
\[
    R^{''} = \begin{pmatrix}
        r'_{11} - b_1 & r'_{12} - b_2 & \dots & r'{1M} - b_m  \\
        \dots \\
        r'_{N1} - b_1 & r'_{N2} - b_2 & \dots & r'{NM} - b_m
        \end{pmatrix}
\]
Ora nota che se  $r'_{11}$ - $b_1$ > 0, e  $r'_{11}$ = $r_{11}$ - $\hat{g}$ \textbf{significa che} 

\[
r_{11} - \hat{g} > b_1,
\]
questo significa che la discossione della recensione dell'utente $1$ sull'item $1$ dalla media globale è maggiore della media $b_j$ , cioè di quanto di discostano dalla media globale le recensioni dei diversi utenti su quel prodotto,
se è maggiore significa che la recensione di quel utente su quel prodotto $j$,è maggiore della media delle recensioni che gli altri utenti hanno assegnato su quello stesso item. 

Questo perchè:
\[
    b_1 = r_{11}-\hat{g} + r_{21}-\hat{g} + \dots + r_{N1} - \hat{g}
\]


\paragraph{Step 5: Calcolo dello User Bias ($b_i$)}
Lo \textbf{User Bias} $b_i$ quantifica la tendenza individuale dell'utente $i$ a fornire valutazioni sistematicamente superiori o inferiori rispetto alla media della community:
\[
b_i = \frac{\sum_{j: (i,j) \in S} r''_{ij}}{N_i + C}
\]
dove $N_i$ rappresenta il numero di rating totali forniti dall'utente $i$.
\begin{itemize}
    \item $b_i > 0$: utente "benevolo" o generoso (assegna voti mediamente alti).
    \item $b_i < 0$: utente "severo" (assegna voti mediamente bassi).
\end{itemize}

\paragraph{Step 6: Calcolo della Matrice Normalizzata / Modello Baseline ($n_{ij}$)}
Unendo le componenti calcolate, si costruisce la stima basale di valutazione $n_{ij}$ (ovvero il modello di stima non personalizzato basato sui soli bias):
\begin{equation}
n_{ij} = \hat{g} + b_j + b_i
\end{equation}

\begin{definition}{Baseline Normalizzata}
La stima baseline $n_{ij}$ esprime il punteggio atteso per la coppia utente-item $(i, j)$ sommando tre componenti additive fondamentali:
\begin{enumerate}
    \item $\hat{g}$: il livello di rating medio dell'intero ecosistema;
    \item $b_j$: la deviazione specifica dovuta al valore intrinseco dell'item $j$;
    \item $b_i$: la deviazione personale legata alla severità o generosità dell'utente $i$.
\end{enumerate}
\end{definition}

Tale modello di baseline costituisce sia un efficace predittore non personalizzato, sia la base di partenza per gli algoritmi personalizzati di \textit{Collaborative Filtering}, i quali operano stimando i residui della matrice sui valori normalizzati $r_{ij} - n_{ij}$.








\section{Algoritmi Personalizzati}

\subsection{Neighborhood-based Collaborative Filtering}
Il Neighborhood-based Collaborative Filtering rappresenta l'approccio classico ai sistemi di raccomandazione di natura collaborativa. Questa metodologia si ramifica in due filoni concettuali:
\begin{itemize}
\item \textbf{User-based:} Il sistema ricerca un sottoinsieme di utenti (il vicinato) con gusti molto affini all'utente target, sfruttandone i rating per suggerire \textbf{nuove alternative}, intuitivamente si confrontano diverse righe della matrice, perchè rappresentano utenti diversi.
\item \textbf{Item-based:} Il sistema ricerca relazioni di somiglianza tra gli item stessi sulla base dei voti ricevuti, utilizzando le preferenze storiche dell'utente per estrapolare un gradimento verso item analoghi.
\end{itemize}
Dal punto di vista operativo, un sistema di Collaborative Filtering può avere due obiettivi principali:

\begin{enumerate}
    \item \textbf{Predizione del rating:} consiste nello stimare il voto che un utente assegnerebbe a uno specifico item, quando tale valutazione non è presente nella User-Rating Matrix (URM) $R$. Il risultato è un singolo valore predetto $\hat{r}_{ij}$, che rappresenta il rating stimato dell'utente $i$ per l'item $j$.
    Si mette il cappelletto ad $r_{ij}$, perchè è il valore stimato, con quello reale che darà l'utente.

    \textit{Esempio:} se Alice ha valutato alcuni film, ma non ha ancora valutato \textit{The Prestige}, il sistema potrebbe stimare che il suo voto per questo film sia $\hat{r}_{ij}=4.8$.

    \item \textbf{Raccomandazione Top-$k$:} consiste nell'individuare e proporre a un utente i $k$ item ritenuti più pertinenti o interessanti. In questo caso, l'obiettivo non è necessariamente prevedere un rating preciso, ma generare una lista ordinata dei migliori $k$ item da raccomandare. Analogamente, è possibile individuare i $k$ utenti ai quali proporre uno specifico item.

    Generalmente questi due tecniche si basano su item, con cui l'utente ancora deve interagire.

    \textit{Esempio:} se il sistema deve raccomandare $k=2$ film ad Alice, potrebbe proporle \textit{The Prestige} e \textit{Dune}, selezionandoli tra quelli con i punteggi di raccomandazione più elevati.

    I due obiettivi sono collegati, ma non equivalenti: un sistema può stimare i rating dei diversi item e poi ordinare i risultati per ottenere una raccomandazione Top-$k$; tuttavia, alcuni algoritmi possono generare direttamente i punteggi di raccomandazione senza predire esplicitamente i rating.
\end{enumerate}

La raccomandazione Top-$k$ è l'approccio più diffuso, poiché lo scopo principale di un \textit{recommender system} è filtrare un insieme potenzialmente molto ampio di item, presentando all'utente soltanto quelli ritenuti più interessanti.

\subsubsection{User-based Neighborhood Collaborative Filtering}
La predizione di un'entry $\hat{r}_{ij}$ (il cui rating effettivo \(r_{ij}\) non è stato osservato) in ambito user-based richiede due step sequenziali:
\begin{enumerate}
\item Trovare l'insieme dei $h$ utenti maggiormente simili all'utente $i$, restrigendo la ricerca solo tra coloro che hanno già fornito un rating all'item $j$.
Cioè trovare gli utenti piu vicini all'utente (cioè che hanno gusti simili) $i$, che hanno già valutato l'item $j$
\item Fondere le valutazioni espresse da questo vicinato per ottenere la stima $\hat{r}_{ij}$.
\end{enumerate}

Per tradurre matematicamente il concetto di affinità, si definisce una misura di similarità $sim(u,v)$ tra un utente $u$ e un utente $v$, basata sugli insiemi $I_u$ e $I_v$ degli item che hanno recensito in comune ($I_u \cap I_v$).
Una prima misura è la \textbf{covarianza} dei rating su questi item condivisi:


$$sim(u,v) = \text{cov}_{I_u \cap I_v}(u,v) = \frac{1}{\vert{}I_u \cap I_v\vert{}} \sum_{k \in I_u \cap I_v} (r_{uk} - \mu_u)(r_{vk} - \mu_v)$$

La covarianza misura se due utenti tendono a dare valutazioni più alte o più basse della propria media agli item in comune.

\begin{itemize}
    \item Se la covarianza è \textbf{positiva} : quando un utente assegna un voto superiore alla propria media, anche l'altro tende a farlo;
    quando un utente assegna un voto inferiore alla propria media, anche l'altro tende a farlo.
    I due utenti tendono a valutare gli stessi item nella stessa direzione rispetto alle proprie medie: se uno apprezza un film più del solito, anche l'altro tende ad apprezzarlo più del solito.
    \item Se la covarianza è \textbf{negativa} : I due utenti tendono a valutare gli stessi item in direzioni opposte rispetto alle proprie medie: quando uno apprezza un film più del solito, l'altro tende ad apprezzarlo meno del solito.
    \item Se la covarianza \textbf{è zero}, non emerge una tendenza lineare sistematica a valutare gli item nella stessa direzione o in direzioni opposte.
\end{itemize}


\begin{center}
\textbf{Questa misura la devi interpretare come "distaccata" dai due punti che abbiamo detto prima, semplicemente questa misura, prende in input due utenti generici}, e considera gli item che hanno valutatato insieme per determinare quanto sono simili, in termini di gusto.
\end{center}

\begin{center}
dove $\mu_u$ e $\mu_v$ indicano i valori medi dei rating degli utenti $u$ e $v$ (sulla riga quindi). Siccome le medie dipendono esclusivamente dal singolo utente e non dalla coppia indagata, si possono calcolare a priori per l'intero sistema, riducendo il carico di calcolo.
\end{center}


La covarianza presenta 3 principali problemi:
\begin{itemize}
    \item \textbf{Aumenta la concordanza reale tra le variabili} : gli scarti \textbf{concordi} dalla media diventano più grandi o più frequenti, quindi sostanzialmente diventa piu frequente lo scarto,
    cioè diventa piu "frequente", che quando un valore è superiore alla propria media, lo è anche l'altro e viceversa.
    \item Si cambia l'unità di misura: passando a un'unità più piccola (es. da metri a centimetri o da chilogrammi a grammi), i numeri crescono e la covarianza si moltiplica di conseguenza.
    \item Aumenta il distacco intorno alla media, cioè i dati diventano piu grandi, aumenta il distacco dalla media, cioè aumenta $(r_{uk} - \mu_u)$ e quindi aumenta tutta la covarianza.
\end{itemize}

Per ovviare a questo problema di scala, si impiega il \textbf{Coefficiente di Correlazione di Pearson}, che normalizza la covarianza introducendo al denominatore le deviazioni standard $\sigma_u$ e $\sigma_v$:

\textbf{La deviazione standard $\sigma$ misura quanto i dati sono dispersi intorno alla media, se la deviazione standard è piccola significa che i dati sono vicini alla media e sono tutti intorno alla media, se la deviazione standard è grande significa che i dati variano tanto}
\paragraph{Esempio}: Immagina due classi di studenti con la stessa media ai test di matematica ($7/10$):
\begin{itemize}
    \item Classe A: i voti sono $6, 7, 7, 7, 8$. I voti sono vicinissimi al $7$. La deviazione standard sarà molto bassa (poca dispersione).
    \item Classe B: i voti sono $4, 5, 7, 9, 10$. Anche qui la media è $7$, ma i voti variano tantissimo. La deviazione standard sarà molto più alta (alta dispersione).
\end{itemize}


$$sim(u,v) = \frac{\text{cov}_{I_u \cap I_v}(u,v)}{\sigma_u \sigma_v} = \frac{\sum_{k \in I_u \cap I_v} (r_{uk} - \mu_u)(r_{vk} - \mu_v)}{\sqrt{\sum_{k \in I_u \cap I_v} (r_{uk} - \mu_u)^2} \sqrt{\sum_{k \in I_u \cap I_v} (r_{vk} - \mu_v)^2}}$$

Il coefficiente di Pearson è compreso tra \(-1\) e \(1\):
\begin{itemize}
    \item +1 : I due utenti seguono esattamente lo stesso andamento nei rating.Indica un andamento perfettamente concorde, non necessariamente voti numericamente uguali, ma significa che quando un utente da un rating sopra la media, lo da anche l'altro.
    \item 0 : Non emerge una relazione lineare tra gli scostamenti dei due utenti.
    \item -1 : Quando un utente si discosta dalla propria media verso l'alto, l'altro si discosta verso il basso, e viceversa.
\end{itemize} 



I risultati di similarità per tutte le coppie di utenti generano la \textit{user-similarity matrix} $\Gamma$, in cui l'elemento $\Gamma_{uv} = sim(u,v)$. La matrice $\Gamma \in \mathbb{R}^{N \times N }$, dove $N$ è il numero degli utenti.
\[
\Gamma =
\begin{pmatrix}
1 & 0.9 & -0.4\\
0.9 & 1 & 0.1\\
-0.4 & 0.1 & 1
\end{pmatrix}
\]
Il coefficiente di $(i,i)$ è 1, perchè un utente ha correlazione massima con sè stesso.





Per limitare l'errore derivante da utenti con metri di giudizio squilibrati, si effettua una predizione centrata sulla media (\textit{mean centered}).
Sia $P_i(j)$ l'insieme dei $k$ utenti più simili a $i$ (\textit{nearest neighbors}) tra coloro che hanno recensito l'item $j$, questo insieme viene calcolato tramite la matrice $\Gamma$, anche se tramite la matrice, non possiamo vedere se uno ha votato l'item $j$ o meno, si vede tramite la matrice originale $R$ gli utenti da selezionare e poi si passa a $\Gamma$.


Supponiamo quindi di voler stimare il rating che Alice assegnerebbe a un item $j$ che non ha ancora valutato. Consideriamo due utenti vicini, Bob e Carlo, (li prendiamo come detto prima dalla matrice $\Gamma$) che hanno già recensito tale item.

Le informazioni a nostra disposizione sono le seguenti:

\begin{center}
\begin{tabular}{|c|c|c|c|}
\hline
\textbf{Utente} & \textbf{Similarità con Alice} & \textbf{Media personale} & \textbf{Rating dell'item $j$} \\
\hline
Bob & $0.9$ & $3$ & $4$ \\
\hline
Carlo & $-0.4$ & $4$ & $3$ \\
\hline
\end{tabular}
\end{center}

\paragraph{1. Calcolo degli scostamenti dalle medie personali}

Per ciascun vicino, calcoliamo la differenza tra il rating assegnato all'item $j$ e la media personale dell'utente. In questo modo, determiniamo se l'utente ha valutato l'item al di sopra o al di sotto del proprio comportamento abituale.

Per Bob, otteniamo:

\[
r_{\text{Bob},j}-\mu_{\text{Bob}}=4-3=+1
\]

Bob ha assegnato un rating pari a $4$, superiore di un punto rispetto alla propria media, pari a $3$.

Per Carlo, invece, otteniamo:

\[
r_{\text{Carlo},j}-\mu_{\text{Carlo}}=3-4=-1
\]

Carlo ha assegnato un rating pari a $3$, inferiore di un punto rispetto alla propria media, pari a $4$.

Questi scostamenti permettono di confrontare i giudizi dei due utenti rispetto alle rispettive abitudini di valutazione, evitando di considerare direttamente i rating assoluti.

\paragraph{2. Pesatura degli scostamenti tramite i coefficienti di Pearson}

Per stimare lo scostamento del rating di Alice, pesiamo gli scostamenti dei vicini mediante i rispettivi coefficienti di similarità:

\[
\Gamma_{\text{Alice,Bob}}=0.9
\]

\[
\Gamma_{\text{Alice,Carlo}}=-0.4
\]

Il contributo di Bob è:

\[
\Gamma_{\text{Alice,Bob}}
(r_{\text{Bob},j}-\mu_{\text{Bob}})
=0.9\cdot(4-3)=0.9
\]

Poiché Bob è positivamente correlato con Alice, il suo scostamento positivo contribuisce positivamente alla predizione.

Il contributo di Carlo è:

\[
\Gamma_{\text{Alice,Carlo}}
(r_{\text{Carlo},j}-\mu_{\text{Carlo}})
=(-0.4)\cdot(3-4)=0.4
\]

In questo caso, sia la similarità sia lo scostamento sono negativi. Il loro prodotto è quindi positivo: la correlazione negativa suggerisce che, quando Carlo valuta un item al di sotto della propria media, Alice tenda a valutarlo al di sopra della sua.

Sommando i contributi, otteniamo il numeratore della formula di predizione:

\[
\sum_{v\in P_i(j)}
\Gamma_{iv}(r_{vj}-\mu_v)
=0.9+0.4=1.3
\]

Entrambi i vicini contribuiscono quindi a una previsione di scostamento positivo per Alice.

\paragraph{3. Normalizzazione tramite la somma delle similarità assolute}

Per ottenere una media pesata, dividiamo la somma dei contributi per la somma dei valori assoluti dei coefficienti di similarità:

\[
\sum_{v\in P_i(j)}|\Gamma_{iv}|
=|0.9|+|-0.4|=0.9+0.4=1.3
\]

Si utilizzano i valori assoluti perché il segno della similarità è già stato considerato nel numeratore. Il denominatore deve invece rappresentare la grandezza complessiva dei pesi, senza che le similarità negative si sottraggano a quelle positive.

La formula completa per la predizione dello scostamento è:

\[
\hat{s}_{ij}=
\frac{
\displaystyle\sum_{v\in P_i(j)}
\Gamma_{iv}(r_{vj}-\mu_v)
}{
\displaystyle\sum_{v\in P_i(j)}|\Gamma_{iv}|
}
\]

Sostituendo i valori del nostro esempio, otteniamo:

\[
\hat{s}_{ij}=\frac{1.3}{1.3}=1
\]

Il risultato indica che il rating di Alice viene stimato superiore di un punto rispetto alla sua media personale.

\paragraph{4. Ricostruzione del rating previsto}

Per ottenere il rating finale, sommiamo lo scostamento stimato alla media dei rating di Alice:

\[
\hat{r}_{ij}=\mu_i+\hat{s}_{ij}
\]

Supponendo che la media di Alice sia $\mu_i=3.5$, otteniamo:

\[
\hat{r}_{ij}=3.5+1=4.5
\]

Pertanto, il sistema prevede che Alice assegnerà all'item $j$ un rating pari a $4.5$.

\begin{center}
    Quindi $\hat{s}_{ij}$ ci dice quanto la predizione dell'utente $i$ sull'item $j$ si discosti dalla media dei suoi voti (quindi dei voti dell'utente $i$) 
\end{center}

\textbf{In sintesi:} il metodo non combina direttamente i rating assoluti dei vicini, ma combina i loro scostamenti dalle rispettive medie, pesandoli in base alla similarità con l'utente target. Infine, aggiunge lo scostamento stimato alla media dell'utente target per ottenere il rating previsto.



Per restituire un voto coerente con la scala assoluta dell'utente $i$, si somma alla fine la sua media di rating originaria:


$$\hat{r}_{ij} = \mu_i + \hat{s}_{ij}$$

\paragraph{Scelta della dimensione del vicinato ($k$)}
La determinazione empirica del parametro $k$ influisce fortemente sulle performance. Fissare un $k$ troppo piccolo riduce drasticamente i dati a disposizione, causando stime volatili. Al contrario, se $k$ è eccessivamente elevato, vengono inclusi nel calcolo utenti con un basso livello di similarità, generando \textit{data noise} che disallinea la previsione. Esiste pertanto un punto di ottimo che massimizza la qualità dell'algoritmo.

\paragraph{Correzioni Avanzate}
Esistono casistiche in cui il calcolo base di similarità restituisce falsi positivi. Per irrobustire la metrica si integrano i seguenti approcci:
\begin{enumerate}
\item \textbf{Discounted Similarity Function:} Una similarità basata su pochissimi item in comune (es. un singolo film) è poco attendibile, anche qualora il coefficiente di Pearson risulti 1. Si definisce quindi una soglia di attendibilità $\beta$; se il numero di item comuni $\vert{}I_u \cap I_v\vert{}$ risulta inferiore a $\beta$, si applica un fattore di sconto proporzionale:


$$\text{DiscountedSim}(u,v) = sim(u,v) \frac{\min(\vert{}I_u \cap I_v\vert{}, \beta)}{\beta}$$


\item \textbf{Popular Item Correction:} Gli item \textit{bestseller} registrano valutazioni ampiamente concordanti tra tutti gli utenti, per cui tenderanno a sovrapporsi frequentemente in $I_u \cap I_v$, mascherando le reali sfumature nei gusti personali[cite: 34]. Assegnando un peso $w_k = \log \left(\frac{m}{m_k}\right)$ (in cui $m$ è il totale dei rating di sistema e $m_k$ è il conteggio dei rating specifici per l'item $k$), è possibile penalizzare gli item troppo popolari direttamente all'interno della formula di Pearson[cite: 34]:
\[sim(u,v) = \frac{\sum_{k \in I_u \cap I_v} w_k (r_{uk} - \mu_u)(r_{vk} - \mu_v)}{\sqrt{\sum_{k \in I_u \cap I_v} (r_{uk} - \mu_u)^2} \sqrt{\sum_{k \in I_u \cap I_v} (r_{vk} - \mu_v)^2}}\]



\end{enumerate}

\subsubsection{Item-based Neighborhood Collaborative Filtering}
Traslando orizzontalmente le metodologie, la valutazione collaborativa può essere applicata sugli item piuttosto che sugli utenti.
Sia $U_h$ l'elenco degli utenti che hanno valutato l'item $h$. Si estrapola la media $\mu_h$ delle recensioni ricevute dal singolo item:


$$\mu_h = \frac{1}{\vert{}U_h\vert{}} \sum_{u \in U_h} r_{uh}$$


L'adattamento del coefficiente di Pearson per misurare l'affinità oggettiva tra l'item $h$ e l'item $k$ diventa:


$$sim(h,k) = \frac{\sum_{u \in U_h \cap U_k} (r_{uh} - \mu_h)(r_{uk} - \mu_k)}{\sqrt{\sum_{u \in U_h \cap U_k} (r_{uh} - \mu_h)^2} \sqrt{\sum_{u \in U_h \cap U_k} (r_{uk} - \mu_k)^2}}$$


Tutti questi valori popolano la \textit{item-similarity matrix} $\Delta_{jk} = sim(j,k)$.
La previsione sfrutta poi l'insieme $Q_j(i)$ dei $k$ item con $\Delta$ maggiore tra quelli valutati dall'utente $i$:


$$\hat{s}_{ij} = \frac{\sum_{k \in Q_j(i)} \Delta_{jk}(r_{ik} - \mu_k)}{\sum_{k \in Q_j(i)} \vert{}\Delta_{jk}\vert{}} \quad \implies \quad \hat{r}_{ij} = \mu_j + \hat{s}_{ij}$$

\subsubsection{Confronto tra User-based e Item-based Collaborative Filtering}
Le architetture commerciali di larga scala presentano un rapporto tra utenti e item enormemente sbilanciato in favore degli utenti. Questo implica che l'intersezione tra chi ha visto l'item $h$ e l'item $k$ sia molto più corposa dell'intersezione tra item visti dall'utente $u$ e $v$ ($\vert{}U_h \cap U_k\vert{} \gg \vert{}I_u \cap I_v\vert{}$).
Da ciò scaturiscono profonde implicazioni tecniche:
\begin{itemize}
\item \textbf{Memory load:} La matrice item-similarity $\Delta$ è di gran lunga meno ingombrante della controparte user-similarity $\Gamma$.
\item \textbf{Stability:} Le entry $\Delta_{jk}$ si fondano su un campione statistico esteso e sono molto resistenti all'inserimento di nuove transazioni, mentre $\Gamma_{uv}$ rischia di fluttuare drasticamente ad ogni interazione. A fronte della crescita degli utenti, $\Gamma$ esige ricalcoli molto più frequenti.
\item \textbf{Accuratezza (Accuracy):} Per lo stesso motivo probabilistico, le raccomandazioni item-based godono tipicamente di una maggiore precisione.
\item \textbf{Explainability:} L'approccio basato sugli item fornisce giustificazioni trasparenti all'utente finale (ad esempio "Visto che hai ascoltato questo brano, ti suggeriamo quest'altro"), risultando esplicativo in modo naturale.
\item \textbf{Novelty e Serendipity:} Unico vero trionfo del modello user-based. Poiché item associati da $\Delta$ tendono a condividere gli stessi archetipi (sequel della stessa saga, film simili), la sorpresa scema. L'approccio user-based scavalca queste limitazioni potendo suggerire un genere radicalmente diverso unicamente perché apprezzato da un "gemello statistico", innalzando la serendipity.
\end{itemize}

\subsection{Modelli a Grafo per Neighborhood Collaborative Filtering}
La bassissima densità dei rating all'interno di una URM inficia l'attendibilità dei classici coefficienti basati sulle preferenze esplicite. Il limite del Pearson risiede nel confrontare esclusivamente nodi con archi diretti condivisi.
Le topologie a grafo $G = (V,E)$ offrono una soluzione strutturale: schematizzando le entità sui nodi $V$ e i pesi sugli archi $E$, è possibile indagare similarità derivate e indirette tra due nodi $h$ e $k$, analizzando i percorsi che fluiscono attraverso molteplici nodi di rimbalzo.

\subsubsection{Item-Item Graphs}
Costruendo un grafo che mappa relazioni tra item, il primo passaggio richiede di legare gli item mediante un calcolo di \textbf{Cosine Similarity} $p'_{h,k}$ tra i vettori colonna $\vec{r}_h$ e $\vec{r}_k$ della URM:


$$p'_{h,k} = \cos(\vec{r}_h, \vec{r}_k) = \frac{\sum_{u \in U} r_{u,h} \cdot r_{u,k}}{\sqrt{\sum_{u \in U} r_{u,h}^2} \cdot \sqrt{\sum_{u \in U} r_{u,k}^2}}$$


Questo genera una matrice di connessione simmetrica. Convertendo tali metriche in transizioni normalizzate, il grafo evolve in una \textbf{Catena di Markov} definita dalla matrice $P$:


$$p_{h,k} = \frac{p'_{h,k}}{\sum_{m \in V} p'_{h,m}}$$


Per distillare da questa catena valori robusti di prossimità vengono esplorate metriche come l'ItemRank e la misura di Katz.

\subsubsection{Misura di Katz}
Importata dall'analisi delle reti, la Misura di Katz calcola l'influenza tra il nodo di partenza $h$ e quello di destinazione $k$ pesando ricorsivamente tutte le vie di collegamento. Si introduce un fattore di smorzamento \textit{damping factor} $\beta \in (0,1)$, affinché le connessioni brevi esercitino un peso esponenzialmente maggiore rispetto a quelle lunghe.
La formulazione classica su una matrice di incidenza $A$ di un grafo orientato è:


$$\Delta_{h,k} = \sum_{l=1}^{\infty} \beta^l \cdot \vert{}paths^l(h,k)\vert{} = \sum_{l=1}^{\infty} \beta^l \cdot A_{h,k}^l$$


Per far convergere la sommatoria, si impone il vincolo $\beta < \frac{1}{\lambda}$, dove $\lambda$ indica il massimo autovalore di $A$.

Riadattando questo concetto matematico all'interno dei sistemi di raccomandazione, si sostituisce la matrice di incidenza con la matrice $P$ della Catena di Markov generata in precedenza. In questa veste, il termine $P_{h,k}^l$ modella la probabilità esatta di transitare da $h$ verso $k$ in esattamente $l$ salti:


$$\Delta_{h,k} = \sum_{l=1}^{\infty} \beta^l \cdot P_{h,k}^l$$


Con questo calcolo esplorativo globale, si ottiene una \textit{item similarity matrix} evoluta e densa, pronta per essere integrata all'interno della classica previsione pesata:


$$\hat{s}_{ij} = \frac{\sum_{k \in Q_j(i)} \Delta_{jk}(r_{ik} - \mu_k)}{\sum_{k \in Q_j(i)} \vert{}\Delta_{jk}\vert{}}$$


dove la selezione dei \textit{nearest item} all'interno del sottoinsieme $Q_j(i)$ avverrà non più per coefficiente locale, bensì in base ai risultati estrapolati dalla Misura di Katz.



\section{Latent Factor Models}

\subsection{Motivazione e idea generale}

Gli approcci di \emph{neighborhood-based collaborative filtering} considerati in precedenza cercano di predire le entry mancanti della \emph{User Rating Matrix} (URM) $R$ utilizzando essenzialmente le valutazioni osservate e le relazioni di similarità tra utenti o tra item. In tali approcci non viene però introdotta, in generale, una particolare ipotesi strutturale sulla matrice $R$.

Questa assenza di struttura può rendere il problema estremamente grande. Indicando con:
\begin{itemize}
    \item $N$ il numero di utenti;
    \item $M$ il numero di item;
    \item $R \in \mathbb{R}^{N \times M}$ la matrice dei rating,
\end{itemize}
il numero complessivo di possibili coppie utente--item, e quindi di entry della matrice $R$, è pari a
\[
N \cdot M.
\]

Come ordine di grandezza, in un contesto come quello di Netflix si può avere:
\[
N \sim 10^8,
\qquad
M \sim 10^3,
\]
da cui
\[
N M \sim 10^{11}.
\]

Il problema originario coinvolge dunque un numero dell'ordine delle centinaia di miliardi di possibili rating.

L'idea fondamentale dei \emph{Latent Factor Models} (LFM) è introdurre una struttura più compatta nella matrice dei rating. Si assume, almeno inizialmente, che $R$ possa essere approssimata mediante il prodotto di due matrici di dimensione molto inferiore:
\[
R \approx XY,
\]
dove
\[
X \in \mathbb{R}^{N \times K},
\qquad
Y \in \mathbb{R}^{K \times M},
\]
e
\[
K \ll N,
\qquad
K \ll M.
\]

Il parametro $K$ rappresenta il numero di \emph{latent factors}, ossia il numero di dimensioni latenti mediante cui si cerca di rappresentare la struttura delle preferenze.

Per una generica entry $(i,j)$ della matrice si ottiene quindi
\begin{equation*}
r_{ij}
\approx
\sum_{k=1}^{K} x_{ik} y_{kj}.
\end{equation*}

In forma vettoriale, indicando con
\[
\mathbf{x}_i
=
(x_{i1},\ldots,x_{iK})
\]
la $i$-esima riga di $X$ e con
\[
\mathbf{y}_j
=
(y_{1j},\ldots,y_{Kj})^{T}
\]
la $j$-esima colonna di $Y$, la predizione può essere scritta come
\[
\widehat r_{ij}
=
\mathbf{x}_i \cdot \mathbf{y}_j.
\]

Pertanto ogni rating viene approssimato attraverso il prodotto scalare tra:
\begin{itemize}
    \item un vettore latente associato all'utente;
    \item un vettore latente associato all'item.
\end{itemize}

\subsubsection{Riduzione della dimensionalità}

Il vantaggio fondamentale della fattorizzazione consiste nella riduzione del numero di parametri.

La matrice originaria $R$ contiene potenzialmente
\[
NM
\]
elementi, mentre le due matrici fattorizzate richiedono complessivamente
\[
NK + KM
=
K(N+M)
\]
parametri.

Riprendendo l'esempio precedente e scegliendo
\[
K=20,
\]
si ottiene
\[
N K + K M
=
10^8 \cdot 20
+
20 \cdot 10^3.
\]

Il numero di parametri diventa quindi dell'ordine dei miliardi, invece che dell'ordine delle centinaia di miliardi. Si ottiene, nell'esempio delle slide, una riduzione di circa due ordini di grandezza rispetto alla dimensione del problema originario.

L'efficienza deriva precisamente dall'ipotesi che le preferenze osservate nella URM non siano completamente indipendenti, ma possano essere spiegate attraverso un numero relativamente piccolo di dimensioni latenti.

\subsection{La matrice predetta e il problema di apprendimento}

Una volta individuate due matrici $X$ e $Y$, è possibile costruire la matrice
\[
\widehat R = XY,
\]
che rappresenta la predizione del modello rispetto alla matrice dei rating reale $R$.

In generale non ci si aspetta che
\[
R = XY.
\]

Al contrario, è altamente improbabile che tutti i rating reali possano essere rappresentati esattamente mediante il prodotto di due matrici di rango ridotto. La fattorizzazione è quindi un'approssimazione:
\[
R \approx \widehat R = XY.
\]

Occorre pertanto introdurre una funzione di errore
\[
\operatorname{Err}(R,\widehat R)
\]
che misuri quanto la matrice predetta $\widehat R$ sia distante dalla matrice reale $R$.

Il problema di apprendimento diventa così un problema di ottimizzazione:
\begin{equation*}
(X^{*},Y^{*})
=
\arg\min_{X,Y}
\operatorname{Err}(R,\widehat R),
\qquad
\widehat R=XY.
\end{equation*}

In altri termini, l'obiettivo è individuare i valori degli elementi di $X$ e $Y$ che consentono di ricostruire nel modo migliore possibile i rating osservati.

\subsection{Intuizione geometrica della fattorizzazione}

La possibilità di esprimere $R$ mediante due matrici più piccole è collegata all'ipotesi che i rating presenti nella URM non siano statisticamente indipendenti.

In termini di algebra lineare, ciò equivale ad assumere che il rango effettivo della matrice $R$ sia inferiore alla massima dimensione teoricamente possibile:
\[
\operatorname{rank}(R)
<
\min\{N,M\}.
\]

La struttura dei dati sarebbe quindi concentrata in uno spazio di dimensione minore rispetto allo spazio originario.

Un'intuizione può essere ottenuta considerando tre film simili, come nell'esempio delle slide:
\begin{itemize}
    \item \emph{Spartacus};
    \item \emph{Gladiator};
    \item \emph{Nero}.
\end{itemize}

Ogni utente può essere rappresentato da un punto in uno spazio tridimensionale:
\[
(r_{\text{Spartacus}},
 r_{\text{Gladiator}},
 r_{\text{Nero}}).
\]

Se i giudizi sui tre film fossero completamente indipendenti, i punti potrebbero occupare liberamente tutto lo spazio tridimensionale. Tuttavia, poiché i tre film sono relativamente simili, è ragionevole aspettarsi delle correlazioni tra i rispettivi rating.

Geometricamente, i punti tenderanno quindi a concentrarsi in prossimità di una struttura di dimensionalità inferiore, ad esempio una retta o un piano nello spazio tridimensionale.

La figura della slide mostra precisamente questa idea: le valutazioni degli utenti sono distribuite prevalentemente lungo una direzione latente. Tale direzione può essere interpretata come un \emph{latent vector} che cattura una componente comune alle valutazioni dei film.

L'informazione presente in più coordinate osservabili può dunque essere sintetizzata attraverso un numero inferiore di coordinate latenti. Questo costituisce l'intuizione geometrica alla base della rappresentazione a rango ridotto.

\subsection{Significato dei fattori latenti}

Partiamo nuovamente dalla relazione
\begin{equation*}
r_{ij}
\approx
\sum_{k=1}^{K} x_{ik}y_{kj}.
\end{equation*}

Essa può essere interpretata come un prodotto scalare:
\[
r_{ij}
\approx
\mathbf{x}_{i}\cdot\mathbf{y}_{j}.
\]

I due vettori sono entrambi costituiti da $K$ componenti.

Si può pensare a ciascuna componente $k$ come a una \emph{feature} o, più precisamente, a un \emph{latent factor}.

In questa interpretazione:
\begin{itemize}
    \item $x_{ik}$ rappresenta il peso o l'importanza che l'utente $i$ attribuisce al fattore latente $k$;
    \item $y_{kj}$ rappresenta quanto il fattore latente $k$ sia presente nell'item $j$.
\end{itemize}

Di conseguenza,
\[
x_{ik}y_{kj}
\]
rappresenta il contributo del fattore $k$ al rating previsto dell'utente $i$ per l'item $j$.

Sommando i contributi di tutti i fattori:
\[
\widehat r_{ij}
=
\sum_{k=1}^{K}x_{ik}y_{kj},
\]
si ottiene la valutazione complessiva prevista.

\subsubsection{Differenza rispetto a un domain model esplicito}

I \emph{Latent Factor Models} presuppongono implicitamente che esistano caratteristiche sottostanti degli item capaci di spiegare le valutazioni degli utenti.

Tuttavia, tali caratteristiche sono profondamente diverse dagli attributi definiti in un tradizionale \emph{domain model}.

In un domain model la struttura viene definita esplicitamente dal progettista. Per esempio, un film potrebbe essere caratterizzato da attributi quali:
\begin{itemize}
    \item genere;
    \item regista;
    \item attori protagonisti;
    \item anno di produzione;
    \item altre caratteristiche note del dominio.
\end{itemize}

Nei Latent Factor Models, invece, i fattori non vengono necessariamente associati a una semantica esplicita. Essi sono innanzitutto astrazioni matematiche apprese dai rating.

Un singolo fattore latente:
\begin{itemize}
    \item potrebbe corrispondere approssimativamente a un attributo interpretabile;
    \item potrebbe non avere alcuna interpretazione immediata;
    \item potrebbe rappresentare una combinazione o mistura di più attributi.
\end{itemize}

La parola \emph{latent} sottolinea precisamente il fatto che tali dimensioni non sono direttamente osservate né imposte a priori.

\subsection{Esempio con due fattori latenti: history e romance}

Le slide introducono un esempio volutamente semplice in cui vengono utilizzati soltanto due fattori latenti:
\[
K=2.
\]

I due fattori sono direttamente interpretabili come generi cinematografici:
\begin{enumerate}
    \item \emph{history};
    \item \emph{romance}.
\end{enumerate}

In questo esempio si utilizza la notazione
\[
R \approx UV^{T},
\]
equivalente alla notazione generale
\[
R \approx XY,
\]
ponendo
\[
X=U,
\qquad
Y=V^{T}.
\]

La matrice $U$ descrive, utente per utente, la relazione con i due generi. Le sue colonne corrispondono quindi ai due fattori latenti \emph{history} e \emph{romance}.

La matrice $V^{T}$ descrive invece quanto ciascun film appartenga ai due generi.

L'esempio considera i film
\begin{itemize}
    \item \emph{Nero};
    \item \emph{Julius Caesar};
    \item \emph{Cleopatra};
    \item \emph{Sleepless in Seattle};
    \item \emph{Pretty Woman};
    \item \emph{Casablanca}.
\end{itemize}

La matrice degli item mostrata nella slide è
\[
V^{T}
=
\begin{pmatrix}
1 & 1 & 1 & 0 & 0 & 0\\
0 & 0 & 1 & 1 & 1 & 1
\end{pmatrix}.
\]

La prima riga rappresenta il fattore \emph{history}, mentre la seconda rappresenta il fattore \emph{romance}.

Pertanto:
\begin{itemize}
    \item \emph{Nero} e \emph{Julius Caesar} sono rappresentati esclusivamente attraverso il fattore \emph{history};
    \item \emph{Sleepless in Seattle}, \emph{Pretty Woman} e \emph{Casablanca} sono rappresentati esclusivamente attraverso il fattore \emph{romance};
    \item \emph{Cleopatra} presenta entrambe le componenti nell'esempio semplificato.
\end{itemize}

La matrice degli utenti riportata nella slide è
\[
U=
\begin{pmatrix}
1  & 0\\
1  & 0\\
1  & 0\\
1  & 1\\
-1 & 1\\
-1 & 1\\
-1 & 1
\end{pmatrix}.
\]

Essa rappresenta tre tipologie schematiche di preferenza:
\begin{itemize}
    \item i primi utenti hanno una preferenza per il fattore \emph{history};
    \item l'utente centrale presenta una preferenza positiva per entrambi i fattori;
    \item gli ultimi utenti presentano una componente negativa rispetto a \emph{history} e positiva rispetto a \emph{romance}.
\end{itemize}

Il prodotto
\[
UV^{T}
\]
produce quindi un'approssimazione della matrice dei rating.

L'esempio mostra anche che tale ricostruzione non deve essere necessariamente perfetta. Definendo la matrice residua
\[
E = R-UV^{T},
\]
la figura successiva mostra che gran parte degli elementi può essere ricostruita esattamente, mentre alcuni rating rimangono differenti dalla previsione.

Nell'esempio grafico, le differenze non nulle sono concentrate nella colonna relativa a \emph{Cleopatra}. Questo evidenzia un punto essenziale: una rappresentazione mediante pochi fattori latenti cattura la struttura dominante nei dati, ma può perdere alcune peculiarità dei singoli rating.

La fattorizzazione è pertanto una rappresentazione approssimata:
\[
R = UV^{T}+E,
\]
dove $E$ raccoglie gli errori di ricostruzione.

\subsection{Scelte progettuali di un Latent Factor Model}

La progettazione di un Latent Factor Model richiede diverse decisioni.

Le slide individuano quattro aspetti principali:
\begin{enumerate}
    \item il numero $K$ dei fattori latenti;
    \item la definizione della funzione di errore;
    \item la tecnica utilizzata per minimizzare l'errore;
    \item eventuali vincoli imposti alle matrici $X$ e $Y$.
\end{enumerate}

Le diverse scelte producono modelli differenti sia per accuratezza sia per tempo di esecuzione.

\subsubsection{Scelta del numero di fattori latenti}

Tipicamente il numero di fattori latenti è dell'ordine
\[
K \approx 10 \div 10^{2}.
\]

La scelta di $K$ determina un compromesso tra capacità espressiva e generalizzazione.

\paragraph{Valore di $K$ troppo piccolo.}

Se $K$ è troppo piccolo, il modello dispone di poche dimensioni latenti e potrebbe non riuscire a rappresentare tutti i fattori che influenzano realmente le valutazioni.

Di conseguenza:
\begin{itemize}
    \item la rappresentazione delle preferenze degli utenti risulta troppo semplice;
    \item utenti con gusti differenti possono essere rappresentati in maniera troppo simile;
    \item il sistema non riesce a personalizzare adeguatamente le raccomandazioni.
\end{itemize}

Il modello presenta quindi una capacità espressiva insufficiente.

\paragraph{Valore di $K$ troppo grande.}

Aumentando $K$, aumenta il numero dei parametri:
\[
NK+KM.
\]

Il modello diventa quindi più flessibile e può adattarsi meglio ai rating osservati.

Tuttavia, soprattutto quando il numero di rating effettivamente osservati è piccolo, un valore troppo elevato di $K$ può generare \emph{overfitting}: il modello riproduce molto bene i dati noti, ma possiede scarso potere predittivo sui rating non osservati.

La scelta di $K$ realizza quindi il classico compromesso tra:
\[
\text{capacità di rappresentazione}
\qquad\text{e}\qquad
\text{capacità di generalizzazione}.
\]

\subsubsection{Possibili condizioni sulle matrici}

A seconda della particolare variante di Matrix Factorization, si possono imporre condizioni differenti sulle matrici $X$ e $Y$.

Le slide citano, a titolo di esempio:
\begin{itemize}
    \item nessuna condizione particolare;
    \item condizioni di ortogonalità;
    \item vincoli di non-negatività:
    \[
    x_{ik}\geq 0,
    \qquad
    y_{kj}\geq 0.
    \]
\end{itemize}

Tali condizioni conducono a differenti varianti del problema di fattorizzazione e possono incidere sia sull'accuratezza sia sul costo computazionale.

\subsection{Funzione di errore}

Una volta stabilito che
\[
\widehat R=XY,
\]
occorre definire quantitativamente la distanza tra la matrice reale $R$ e quella ricostruita $\widehat R$.

Una scelta naturale è utilizzare la norma di Frobenius al quadrato.

Se tutte le entry di $R$ fossero note, si potrebbe definire
\begin{equation*}
\operatorname{Err}(R,\widehat R)
=
\frac{1}{2}
\lVert R-\widehat R\rVert_F^2
=
\frac{1}{2}
\lVert R-XY\rVert_F^2.
\end{equation*}

La norma di Frobenius soddisfa
\[
\lVert A\rVert_F^2
=
\sum_i\sum_j a_{ij}^{2},
\]
per cui essa corrisponde alla somma dei quadrati degli scarti tra le entry delle due matrici.

Il fattore
\[
\frac{1}{2}
\]
non modifica il punto di minimo della funzione, ma rende più semplici le espressioni delle derivate, poiché la derivata del termine quadratico cancella il fattore $2$.

\subsubsection{URM incompleta}

Nella pratica, però, la URM è sparsa: soltanto una piccola parte dei rating è osservata.

Si definisce quindi l'insieme
\[
S
=
\{(i,j)\mid r_{ij}\text{ è osservato}\}.
\]

La funzione di errore viene calcolata esclusivamente sulle coppie appartenenti a $S$.

Definiamo innanzitutto l'errore relativo alla singola entry:
\[
e_{ij}
=
r_{ij}-\widehat r_{ij}.
\]

Poiché
\[
\widehat r_{ij}
=
\sum_{k=1}^{K}x_{ik}y_{kj},
\]
si ha
\[
e_{ij}
=
r_{ij}
-
\sum_{k=1}^{K}x_{ik}y_{kj}.
\]

La funzione di errore diventa quindi
\begin{align*}
\operatorname{Err}(R,\widehat R)
&=
\frac{1}{2}
\sum_{(i,j)\in S}e_{ij}^{2}\\
&=
\frac{1}{2}
\sum_{(i,j)\in S}
\left(r_{ij}-\widehat r_{ij}\right)^2\\
&=
\frac{1}{2}
\sum_{(i,j)\in S}
\left(
r_{ij}
-
\sum_{k=1}^{K}x_{ik}y_{kj}
\right)^2.
\end{align*}

La differenza rispetto alla norma di Frobenius completa è fondamentale: le entry non osservate non contribuiscono direttamente alla funzione obiettivo.

\subsection{Missing as Random Assumption}

Il fatto di calcolare l'errore soltanto sull'insieme
\[
S
\]
implica un'assunzione statistica importante.

Si assume che i rating osservati costituiscano un campione sufficientemente rappresentativo dell'insieme complessivo dei rating.

In altri termini, si ipotizza che i rating non osservati abbiano, dal punto di vista statistico, un comportamento abbastanza simile a quelli osservati. Se fosse possibile inserirli nel calcolo dell'errore, essi non dovrebbero modificare significativamente la stima complessiva.

Questa ipotesi è indicata nelle slide come \emph{missing as random assumption}.

L'assunzione può tuttavia essere problematica in sistemi reali.

Nel dominio dei film, per esempio:
\begin{itemize}
    \item un utente può essere maggiormente portato a recensire i film che gli sono piaciuti;
    \item alcuni film possono essere presentati all'utente ma successivamente scartati;
    \item un film scartato potrebbe presumibilmente ricevere una valutazione bassa, pur non comparendo nella URM come rating esplicito;
    \item di conseguenza, ci si può aspettare che il rating medio dei film non recensiti sia leggermente inferiore al rating medio dei film recensiti.
\end{itemize}

Se la differenza tra la distribuzione dei rating osservati e quella dei rating mancanti è piccola, l'ipotesi \emph{missing as random} può costituire comunque una buona approssimazione.

Se invece il processo che genera i rating mancanti introduce un bias rilevante, l'errore calcolato esclusivamente su $S$ può non rappresentare adeguatamente l'errore reale sull'intera matrice.

\subsection{Regularization}

La minimizzazione dell'errore sui soli rating osservati può essere interpretata come un problema di interpolazione: si cerca di adattare il prodotto
\[
XY
\]
ai valori noti della matrice $R$.

Se la URM è molto sparsa e il numero di fattori $K$ è sufficientemente elevato, il modello può disporre di un numero di gradi di libertà tale da adattarsi eccessivamente ai rating appartenenti a $S$.

In particolare, quando
\[
|S| \ll NM,
\]
il rischio di \emph{overfitting} è significativo.

Per contrastarlo si introducono nella funzione obiettivo dei termini di \emph{regularization}, che penalizzano soluzioni eccessivamente adattate ai dati osservati.

La funzione regolarizzata riportata nelle slide è
\begin{equation*}
\operatorname{Err}(R,\widehat R)
=
\lVert R-XY\rVert_{S}^{2}
+
\lambda_{1}\lVert X\rVert_{F}^{2}
+
\lambda_{2}\lVert Y\rVert_{F}^{2},
\end{equation*}
dove la notazione
\[
\lVert R-XY\rVert_{S}^{2}
\]
indica che l'errore di ricostruzione viene valutato esclusivamente sulle entry osservate appartenenti a $S$.

Esplicitando i termini, con la convenzione del fattore $1/2$, si ottiene
\begin{align*}
\operatorname{Err}(R,\widehat R)
={}&
\frac{1}{2}
\sum_{(i,j)\in S}e_{ij}^{2}\\
&+
\frac{\lambda_{1}}{2}
\sum_{i=1}^{N}\sum_{k=1}^{K}x_{ik}^{2}\\
&+
\frac{\lambda_{2}}{2}
\sum_{k=1}^{K}\sum_{j=1}^{M}y_{kj}^{2}.
\end{align*}

Equivalentemente:
\[
J(X,Y)
=
\frac{1}{2}
\sum_{(i,j)\in S}
\left(
r_{ij}
-
\sum_{k=1}^{K}x_{ik}y_{kj}
\right)^2
+
\frac{\lambda_1}{2}\lVert X\rVert_F^2
+
\frac{\lambda_2}{2}\lVert Y\rVert_F^2.
\]

I coefficienti
\[
\lambda_1\geq 0,
\qquad
\lambda_2\geq 0
\]
controllano l'intensità della regolarizzazione.

Quanto maggiori sono $\lambda_1$ e $\lambda_2$, tanto più vengono favorite soluzioni caratterizzate da valori piccoli degli elementi
\[
x_{ik}
\qquad\text{e}\qquad
y_{kj}.
\]

La regularization introduce quindi un compromesso tra:
\begin{itemize}
    \item adattamento ai rating osservati;
    \item complessità del modello.
\end{itemize}

Con coefficienti di regolarizzazione piccoli, il modello è libero di adattarsi maggiormente ai dati di training. Con coefficienti più grandi viene invece penalizzata una crescita eccessiva dei parametri latenti.

\subsection{Minimizzazione mediante Gradient Descent}

Per individuare le matrici $X$ e $Y$ che minimizzano la funzione obiettivo vengono utilizzate varianti del metodo del \emph{Gradient Descent}.

Per semplicità, indichiamo la funzione di errore mediante
\[
J
=
\operatorname{Err}(R,\widehat R).
\]

Ricordiamo che
\[
e_{ij}
=
r_{ij}
-
\sum_{k=1}^{K}x_{ik}y_{kj}.
\]

\subsubsection{Derivata rispetto a un fattore dell'utente}

Consideriamo un generico parametro
\[
x_{iq},
\]
ossia la componente $q$ del vettore latente relativo all'utente $i$.

Senza regularization:
\[
J
=
\frac{1}{2}
\sum_{(i,j)\in S}e_{ij}^{2}.
\]

Soltanto i rating che coinvolgono l'utente $i$ dipendono da $x_{iq}$. Per tali termini:
\[
e_{ij}
=
r_{ij}
-
\sum_{k=1}^{K}x_{ik}y_{kj}.
\]

La derivata dell'errore rispetto a $x_{iq}$ è
\[
\frac{\partial e_{ij}}{\partial x_{iq}}
=
-y_{qj}.
\]

Applicando la regola della catena:
\begin{align*}
\frac{\partial}{\partial x_{iq}}
\left(
\frac{1}{2}e_{ij}^{2}
\right)
&=
e_{ij}
\frac{\partial e_{ij}}{\partial x_{iq}}\\
&=
-y_{qj}e_{ij}.
\end{align*}

Sommando su tutti gli item $j$ valutati dall'utente $i$:
\begin{equation*}
\frac{\partial J}{\partial x_{iq}}
=
-
\sum_{j:(i,j)\in S}
y_{qj}e_{ij}.
\end{equation*}

Introducendo il termine di regularization:
\[
\frac{\lambda_1}{2}\sum_{i,k}x_{ik}^{2},
\]
la sua derivata rispetto a $x_{iq}$ è
\[
\lambda_1x_{iq}.
\]

Pertanto:
\begin{equation*}
\boxed{
\frac{\partial J}{\partial x_{iq}}
=
-
\sum_{j:(i,j)\in S}
y_{qj}e_{ij}
+
\lambda_1x_{iq}
}
\end{equation*}

\subsubsection{Derivata rispetto a un fattore dell'item}

Consideriamo ora
\[
y_{qj},
\]
ossia la componente $q$ del vettore latente associato all'item $j$.

Analogamente:
\[
\frac{\partial e_{ij}}{\partial y_{qj}}
=
-x_{iq}.
\]

In assenza di regularization:
\begin{equation*}
\frac{\partial J}{\partial y_{qj}}
=
-
\sum_{i:(i,j)\in S}
x_{iq}e_{ij}.
\end{equation*}

La regularization su $Y$ introduce il termine
\[
\frac{\lambda_2}{2}
\sum_{k,j}y_{kj}^{2},
\]
la cui derivata rispetto a $y_{qj}$ è
\[
\lambda_2y_{qj}.
\]

Quindi:
\begin{equation*}
\boxed{
\frac{\partial J}{\partial y_{qj}}
=
-
\sum_{i:(i,j)\in S}
x_{iq}e_{ij}
+
\lambda_2y_{qj}
}
\end{equation*}

Le formule mettono in evidenza che:
\begin{itemize}
    \item l'aggiornamento di un vettore utente dipende soltanto dai rating osservati di quell'utente;
    \item l'aggiornamento di un vettore item dipende soltanto dai rating osservati relativi a quell'item.
\end{itemize}

\subsubsection{Matrici gradiente}

Si definisce
\[
\nabla J_X
\]
come la matrice $N\times K$ contenente tutte le derivate
\[
\frac{\partial J}{\partial x_{iq}},
\qquad
i\in\{1,\ldots,N\},
\quad
q\in\{1,\ldots,K\}.
\]

In forma esplicita:
\[
\nabla J_X
=
\left[
\frac{\partial J}{\partial x_{iq}}
\right]_
{\substack{i=1,\ldots,N\\q=1,\ldots,K}}.
\]

Analogamente,
\[
\nabla J_Y
\]
è la matrice $K\times M$ delle derivate
\[
\frac{\partial J}{\partial y_{qj}},
\qquad
q\in\{1,\ldots,K\},
\quad
j\in\{1,\ldots,M\}.
\]

Queste due matrici descrivono la direzione di massima crescita della funzione $J$ nello spazio dei parametri. Per minimizzare $J$ occorre pertanto muoversi nella direzione opposta al gradiente.

\subsection{Procedura generale di Gradient Descent}

In linea generale, l'apprendimento mediante gradiente discendente segue i seguenti passaggi.

\begin{enumerate}
    \item \textbf{Inizializzazione.}

    Si assegnano valori iniziali alle matrici
    \[
    X
    \qquad\text{e}\qquad
    Y.
    \]

    Le slide indicano, come possibilità, un'inizializzazione casuale.

    \item \textbf{Calcolo dei gradienti.}

    A partire dai valori correnti di $X$ e $Y$, si calcolano
    \[
    \nabla J_X
    \qquad\text{e}\qquad
    \nabla J_Y.
    \]

    \item \textbf{Aggiornamento dei parametri.}

    Le matrici vengono aggiornate nella direzione opposta rispetto al gradiente:
    \begin{align*}
    X &\leftarrow X-\alpha\nabla J_X,\\
    Y &\leftarrow Y-\alpha\nabla J_Y.
    \end{align*}

    \item \textbf{Iterazione.}

    Il processo viene ripetuto utilizzando i nuovi valori delle matrici.

    \item \textbf{Terminazione.}

    L'algoritmo termina quando viene soddisfatta una opportuna condizione di convergenza.
\end{enumerate}

Il parametro
\[
\alpha>0
\]
è detto \emph{step size}.

Esso controlla l'ampiezza dell'aggiornamento effettuato a ogni iterazione.

Le slide sottolineano che il valore preciso di $\alpha$ dipende dalla particolare tecnica utilizzata, ma è generalmente piccolo; viene riportato come esempio:
\[
\alpha=0.005.
\]

Un singolo aggiornamento di un parametro può essere interpretato schematicamente come
\[
\theta
\leftarrow
\theta
-
\alpha\frac{\partial J}{\partial\theta},
\]
dove $\theta$ rappresenta un generico elemento di $X$ o di $Y$.

\subsection{Principali strategie di minimizzazione}

Le diverse tecniche di apprendimento si distinguono soprattutto per il modo in cui vengono aggiornate le matrici $X$ e $Y$.

Le slide introducono tre principali famiglie:
\begin{enumerate}
    \item aggiornamento \emph{Batch};
    \item aggiornamento \emph{Stochastic};
    \item \emph{Alternating Least Squares} (ALS).
\end{enumerate}

\subsubsection{Batch Gradient Descent}

Nell'approccio \emph{Batch} il gradiente viene calcolato utilizzando globalmente l'insieme dei rating osservati.

Concettualmente, si calcolano
\[
\nabla J_X
\qquad\text{e}\qquad
\nabla J_Y
\]
sulla base dell'intero insieme
\[
S
\]
prima di effettuare il corrispondente aggiornamento.

L'aggiornamento generale è:
\begin{align*}
X &\leftarrow X-\alpha\nabla J_X,\\
Y &\leftarrow Y-\alpha\nabla J_Y.
\end{align*}

In tale strategia, ogni passo di apprendimento riflette quindi l'informazione proveniente da tutte le entry osservate.

\subsubsection{Stochastic Gradient Descent}

Nella variante stocastica, invece di calcolare a ogni passo le sommatorie su tutto $S$, viene selezionata una coppia osservata
\[
(i,j)\in S,
\]
tipicamente in maniera casuale.

Sulla base del solo rating $r_{ij}$ vengono aggiornati i parametri interessati:
\[
(x_{i1},\ldots,x_{iK}),
\]
ossia il vettore latente dell'utente $i$, e
\[
(y_{1j},\ldots,y_{Kj}),
\]
ossia il vettore latente dell'item $j$.

La tecnica è detta \emph{stochastic} proprio perché l'aggiornamento viene effettuato utilizzando una singola osservazione selezionata dal dataset invece dell'intero insieme dei rating.

Le derivate complete
\[
-
\sum_{j:(i,j)\in S}y_{qj}e_{ij}
\qquad\text{e}\qquad
-
\sum_{i:(i,j)\in S}x_{iq}e_{ij}
\]
vengono quindi approssimate attraverso contributi locali relativi alla coppia estratta.

\subsubsection{Alternating Least Squares}

L'\emph{Alternating Least Squares} (ALS) segue una logica differente.

Il problema
\[
R\approx XY
\]
dipende simultaneamente da due matrici incognite. ALS procede alternando l'ottimizzazione delle due matrici:
\begin{enumerate}
    \item si fissa $X$ e si apprende $Y$;
    \item si fissa $Y$ e si apprende $X$;
    \item il processo viene ripetuto fino alla convergenza.
\end{enumerate}

In forma schematica:
\[
\boxed{
X\ \text{fissata}
\quad\Longrightarrow\quad
\text{apprendimento di }Y
}
\]
e successivamente
\[
\boxed{
Y\ \text{fissata}
\quad\Longrightarrow\quad
\text{apprendimento di }X.
}
\]

L'aspetto caratterizzante è dunque l'alternanza: anziché modificare contemporaneamente entrambe le matrici, si ottimizza una matrice mantenendo temporaneamente costante l'altra.

\subsection{Sintesi del modello}

Un Latent Factor Model costruisce una rappresentazione a bassa dimensionalità della User Rating Matrix:
\[
R\approx XY,
\]
con
\[
X\in\mathbb{R}^{N\times K},
\qquad
Y\in\mathbb{R}^{K\times M},
\qquad
K\ll\min\{N,M\}.
\]

La predizione di un rating è
\[
\widehat r_{ij}
=
\sum_{k=1}^{K}x_{ik}y_{kj}
=
\mathbf{x}_i\cdot\mathbf{y}_j.
\]

L'apprendimento consiste nel determinare $X$ e $Y$ minimizzando una funzione obiettivo definita sui rating osservati:
\[
J(X,Y)
=
\frac{1}{2}
\sum_{(i,j)\in S}
\left(
r_{ij}
-
\sum_{k=1}^{K}x_{ik}y_{kj}
\right)^2,
\]
eventualmente aggiungendo termini di regularization:
\[
J(X,Y)
=
\frac{1}{2}
\sum_{(i,j)\in S}e_{ij}^{2}
+
\frac{\lambda_1}{2}\lVert X\rVert_F^2
+
\frac{\lambda_2}{2}\lVert Y\rVert_F^2.
\]

I gradienti fondamentali sono
\[
\frac{\partial J}{\partial x_{iq}}
=
-
\sum_{j:(i,j)\in S}
y_{qj}e_{ij}
+
\lambda_1x_{iq},
\]
e
\[
\frac{\partial J}{\partial y_{qj}}
=
-
\sum_{i:(i,j)\in S}
x_{iq}e_{ij}
+
\lambda_2y_{qj}.
\]

L'aggiornamento mediante Gradient Descent assume quindi la forma
\begin{align*}
X &\leftarrow X-\alpha\nabla J_X,\\
Y &\leftarrow Y-\alpha\nabla J_Y.
\end{align*}

Le principali scelte progettuali riguardano:
\begin{itemize}
    \item il numero di fattori latenti $K$;
    \item la funzione di errore;
    \item l'intensità della regularization;
    \item gli eventuali vincoli su $X$ e $Y$;
    \item la strategia di minimizzazione, ad esempio Batch, Stochastic Gradient Descent o ALS.
\end{itemize}

L'idea complessiva è che una matrice di rating apparentemente enorme e sparsa possa possedere una struttura intrinseca molto più semplice: le preferenze degli utenti e le caratteristiche degli item vengono rappresentate nello stesso spazio latente di dimensione $K$, e la loro compatibilità viene stimata attraverso il prodotto scalare dei rispettivi vettori.